\documentclass[11pt]{article}

% --- Encoding / fonts ---
\usepackage[T1]{fontenc}
\usepackage[utf8]{inputenc}
\usepackage{lmodern}
\usepackage{microtype}

% --- Math ---
\usepackage{amsmath,amssymb,amsthm,mathtools}
\usepackage{bm}

% --- Page layout ---
\usepackage[margin=1in]{geometry}
\usepackage{parskip}
\setlength{\parskip}{0.5\baselineskip plus 2pt minus 1pt}
\setlength{\parindent}{0pt}

% --- Tables ---
\usepackage{booktabs}
\usepackage{longtable}
\usepackage{array}

% --- Lists ---
\providecommand{\tightlist}{\setlength{\itemsep}{0pt}\setlength{\parskip}{0pt}}

% --- Hyperref (must be loaded near the end) ---
\usepackage{xcolor}
\usepackage[hidelinks,unicode,linktoc=all]{hyperref}

% --- Numbering depth: section / subsection / subsubsection ---
\setcounter{secnumdepth}{3}
\setcounter{tocdepth}{2}

% --- Better widow/orphan control ---
\widowpenalty=10000
\clubpenalty=10000

% --- Title metadata ---
\title{\Large\bfseries CD-SBI: Calibrated Confidence Distributions \\ via Simulation-Based Inference}
\author{}
\date{\normalsize\itshape A working technical document. Comments, counter-examples, and pushback welcome.}

\begin{document}

\maketitle

\thispagestyle{empty}
\vspace{1em}

\section*{Abstract}

Simulation-based inference (SBI) is now the standard tool for
likelihood-free inference in scientific simulators, but its commonly
used objectives --- neural posterior, likelihood, and ratio estimation
--- do not enforce frequentist coverage. Hermans et al.~(2022)
documented that off-the-shelf SBI methods are often substantially
overconfident at realistic simulation budgets, with no internal
diagnostic to detect this.

This document develops an alternative framework whose training target
is, by construction, a calibrated \emph{confidence distribution} (CD) in
the Schweder--Hjort / Fraser tradition. The method learns a pivotal
quantity \(r: \Theta \times \mathcal{X} \to \mathbb{R}^d\) such that,
under the data-generating distribution,
\(r(\theta_0; X) \mid \theta_0 \sim \mathcal{N}(0, I_d)\) for every true
\(\theta_0\). The induced CD then has exact frequentist coverage at
every confidence level --- \emph{pointwise in \(\theta_0\), not merely
on average over the proposal.}

The method has three ingredients:

\begin{enumerate}
\def\labelenumi{\arabic{enumi}.}
\tightlist
\item
  \textbf{A standard NF-MLE objective}: per-sample maximum likelihood
  for \(r\) viewed as a conditional normalizing flow,
  \[L(r) \;=\; \mathbb{E}_{(\theta, X) \sim \rho \otimes P_\theta}\!\left[\tfrac{1}{2}\|r(\theta, X)\|^2 \;-\; \log\bigl|\det \partial r/\partial X\bigr|\right].\]
  This is the same objective used by Sequential Neural Likelihood (SNL;
  Papamakarios et al.~2019).
\item
  \textbf{Architectural monotonicity in \(\theta\)}: \(r\) is
  parameterized so that \(\theta \mapsto r_k(\theta; X)\) is strictly
  increasing in each coordinate (autoregressively, in higher
  dimensions). This is what converts a trained flow into a valid CD
  pivot.
\item
  \textbf{Architectural monotonicity in \(X\)}: \(r\) is parameterized
  so that \(\partial r/\partial X\) is (lower-)triangular with positive
  diagonal. This is what makes the change-of-variables Jacobian both
  correct and computationally cheap, and avoids the spurious sub-optima
  that arise when monotonicity-in-\(X\) is left to autograd.
\end{enumerate}

We establish that:

\begin{itemize}
\tightlist
\item
  In the regular continuous one-parameter case under \(C^1\) regularity,
  the NF-MLE objective combined with monotonicity-in-\(\theta\) uniquely
  recovers the Schweder--Hjort UMP-unbiased CD (Theorem A). A Lipschitz
  version handles ReLU networks under one mild side condition (Theorem
  A*).
\item
  The result lifts to any one-parameter regular exponential family with
  continuous sufficient statistic (Theorem C), and to discrete families
  via Lancaster's randomized PIT (Theorem C*).
\item
  In multivariate \(\theta\), a triangular-autoregressive construction
  gives a tractable architecture whose Jacobian factors into a product
  of \(d\) positive scalars, scaling linearly in \(d\). Within that
  architectural class, the calibration manifold reduces to the unique
  Knothe--Rosenblatt rearrangement (Theorem A-d).
\item
  Empirically, the method recovers the true pivot to within
  \textasciitilde1--3\% relative error on 1D and 2D location-Gaussian
  problems and on the exponential-rate model, with empirical coverage
  matching nominal levels to within 0.01--0.02 across all tested
  parameter values and confidence levels.
\end{itemize}

The framework reduces to the classical optimal CD in regular cases and
admits a tractable amortized neural implementation in cases where no
closed-form pivot is available. We document one architectural failure
mode (the autograd-Jacobian shortcut without monotone-in-\(X\)
enforcement) and its fix.

\medskip\hrule\medskip

\tableofcontents

\medskip\hrule\medskip

\section*{Part I — Framework}
\addcontentsline{toc}{section}{Part I — Framework}

\section{Motivation: SBI and confidence distributions}\label{sec:1}

\subsection{The state of SBI}\label{subsec:1.1}

Many scientific models are \emph{implicit}: defined by a stochastic
simulator \(X \sim P_\theta\) from which one can sample data, but for
which the likelihood \(p(X \mid \theta)\) is intractable or unavailable
in closed form. Modern \emph{simulation-based inference} (SBI) trains a
deep-network surrogate for the quantity of interest --- the Bayesian
posterior \(p(\theta \mid X)\), the likelihood \(p(X \mid \theta)\), or
the likelihood ratio --- from simulated pairs \((\theta, X)\).

The three dominant families are:

\begin{itemize}
\tightlist
\item
  \textbf{Neural posterior estimation} (NPE; Papamakarios \& Murray
  2016; Greenberg et al.~2019): train a density estimator directly on
  the posterior \(\theta \mid X\).
\item
  \textbf{Neural likelihood estimation / Sequential Neural Likelihood}
  (NLE/SNL; Papamakarios et al.~2019): train a density estimator on
  \(X \mid \theta\), then run MCMC over the learned surrogate.
\item
  \textbf{Neural ratio estimation} (NRE; Hermans et al.~2020; Miller et
  al.~2022): train a binary classifier to distinguish
  \((\theta, X) \sim p(\theta, X)\) from
  \((\theta, X) \sim p(\theta)p(X)\), recovering the likelihood ratio
  via Bayes's rule.
\end{itemize}

On benchmarks with adequate simulation budgets, these methods are
accurate.

The problem is \emph{miscalibration}. Hermans et al.~(2022) (``A Trust
Crisis in Simulation-Based Inference?'') documented that all four major
SBI families systematically produce \emph{overconfident} posteriors at
realistic budgets --- credible intervals have empirical coverage well
below nominal. Subsequent work has tried to retrofit calibration via
conservative regularization (Delaunoy et al.~2022), differentiable
coverage probabilities (Falkiewicz et al.~2023), Neyman-construction
wrappers (LF2I: Dalmasso et al.~2024; WALDO: Masserano et al.~2023),
depth-based hyper-rectangles (Bortolato \& Ventura 2025), and
variational coverage corrections (Patel et al.~2023). These methods all
treat calibration as a correction to a non-calibrated primary objective.

\subsection{A different starting point: confidence distributions}\label{subsec:1.2}

We instead change the primary target of inference. Rather than
estimating a posterior and correcting it, we aim to produce --- from the
simulator alone --- a \emph{confidence distribution}.

\textbf{What a confidence distribution is.} A confidence distribution
(CD) is the frequentist analogue of the Bayesian posterior: a
data-dependent distribution function \(H(\cdot; X)\) on \(\Theta\) whose
\(\alpha\)-quantile is the upper endpoint of an exact one-sided
\(\alpha\)-confidence interval. The modern theory traces back through
Fisher's fiducial inference (1930), Cox (1958), and Efron (1993, 1998),
to the synthesis in Schweder \& Hjort (2016). Schweder--Hjort
(2016, Def.~3.1) require two properties: (i) \(H(\cdot; X)\) is a CDF
on \(\Theta\) for each \(X\), and (ii) the calibration condition
\textbf{exact coverage}: \[
H(\theta_0; X) \;\sim\; U(0,1) \quad \text{when } X \sim P_{\theta_0}, \text{ for every } \theta_0 \in \Theta.
\] Property (i) is enforced architecturally in this paper by (R1)
(\S\ref{subsec:2.2}); property (ii) is what NF-MLE drives the network
toward. Equivalently to (ii): \(H(\theta_0; X)\) is the
\emph{probability integral transform} (PIT) --- the mapping that
converts any continuous random variable into a uniform \((0,1)\)
variable via its own CDF --- applied to a correctly chosen
\emph{pivot}.

\textbf{Astronomer-friendly comparison.} A CD plays the same role as a
posterior: a distribution-valued summary of the evidence for \(\theta\).
The difference is what calibration means.

\begin{itemize}
\tightlist
\item
  A Bayesian posterior \(\pi(\theta \mid X)\) depends on a prior. Its
  credible intervals contain the true \(\theta_0\) with probability
  \(\alpha\) \emph{averaged over the prior} --- coverage holds when one
  integrates over the same prior used to define the posterior. There is
  generally no guarantee of frequentist coverage \emph{pointwise at a
  given \(\theta_0\)}.
\item
  A CD \(H(\theta; X)\) depends only on the model \(\{P_\theta\}\) and
  the data. Its quantile-based intervals contain \(\theta_0\) with
  probability \emph{exactly} \(\alpha\) \emph{for every fixed
  \(\theta_0\)}, the data being random.
\end{itemize}

In regular one-parameter exponential families with monotone likelihood
ratio in a continuous sufficient statistic, the optimal (uniformly most
powerful unbiased, ``UMPU'') CD is explicit and unique (the precise
hypotheses are stated in \S\ref{subsec:5.1}; the location-normal model
is the simplest case). For \(X \sim \mathcal{N}(\theta, 1)\) it is \[
H^*(\theta; X) = \Phi(\theta - X),
\] where \(\Phi\) is the standard normal CDF. Check: under
\(X \sim \mathcal{N}(\theta_0, 1)\) one has \(\theta_0 - X = -Z\) with
\(Z \sim \mathcal{N}(0,1)\), so
\(\Phi(\theta_0 - X) = \Phi(-Z) \sim U(0,1)\) exactly. Inverting at
level \(1-\alpha/2\) and \(\alpha/2\) recovers the textbook
\(z\)-interval \([X - z_{1-\alpha/2}, X + z_{1-\alpha/2}]\).

Our aim is to compute such a CD from a simulator alone, when no closed
form is available.

\subsection{The top-line goal and what supports it}\label{subsec:1.3}

\textbf{Top-line goal.} Produce, from a simulator alone, a function
\(r(\theta; X)\) whose induced CD \[
H_r(\theta; X) \;:=\; \Phi_d(r(\theta; X))
\] (\(\Phi_d\) = product standard normal CDF) has \emph{pointwise exact
frequentist coverage} --- every confidence level, every true
\(\theta_0\).

The object \(r\) is called a \emph{pivot}: in classical statistics, a
function of data and parameter whose distribution does not depend on the
parameter (the canonical example being \(X - \theta\) in the
location-normal model). Here we require the slightly stronger property
that its distribution under the truth is exactly
\(\mathcal{N}(0, I_d)\).

\textbf{Two characterizing properties.}

\begin{itemize}
\tightlist
\item
  \textbf{(C1) Calibration.} Under \(X \sim P_{\theta_0}\),
  \(r(\theta_0; X) \sim \mathcal{N}(0, I_d)\) for every
  \(\theta_0 \in \Theta\). Equivalently,
  \(H_r(\theta_0; X) \sim U(0,1)^d\). This is a \emph{frequentist}
  property --- no prior, no expectation over \(\theta\).
\item
  \textbf{(C2) Monotonicity in \(\theta\).} In 1D,
  \(\theta \mapsto r(\theta; X)\) is strictly increasing. In \(d > 1\),
  \(\theta_k \mapsto r_k(\theta; X)\) is strictly increasing at fixed
  \(\theta_{<k}\) and \(X\) (the autoregressive form developed in
  \S\ref{sec:6}). The autoregressive form is what makes the joint
  \(\Phi_d(r(\cdot; X))\) a Knothe--Rosenblatt-type transport with
  connected level sets; coordinate-wise monotonicity alone (with all
  other \(\theta_j\) held fixed) is insufficient in \(d > 1\).
\end{itemize}

\textbf{Subordinate objectives the rest of the document develops in
service of this goal.}

\begin{longtable}[]{@{}
  >{\raggedright\arraybackslash}p{(\columnwidth - 2\tabcolsep) * \real{0.5000}}
  >{\raggedright\arraybackslash}p{(\columnwidth - 2\tabcolsep) * \real{0.5000}}@{}}
\toprule\noalign{}
\begin{minipage}[b]{\linewidth}\raggedright
Objective
\end{minipage} & \begin{minipage}[b]{\linewidth}\raggedright
Where addressed
\end{minipage} \\
\midrule\noalign{}
\endhead
\bottomrule\noalign{}
\endlastfoot
Characterize the space of \(r\)'s satisfying (C1): the \emph{calibration
manifold} \(\mathcal{M}\) & \S\ref{sec:2} \\
Identify a loss whose minimizers on the relevant function class are
\(\mathcal{M}\): NF-MLE & \S\ref{sec:3}, \S\ref{subsec:3.7} \\
Identify architectural constraints picking out a unique element of
\(\mathcal{M}\) & \S\ref{subsec:3.3}, \S\ref{subsec:3.5}, \S\ref{sec:6} \\
Prove uniqueness in 1D regular models & \S\ref{sec:4}--\S\ref{sec:5} (Theorems A, A\emph{, C,
C}) \\
Extend to multivariate \(\theta\) & \S\ref{sec:6} (Theorem A-d via
Knothe--Rosenblatt) \\
Validate at finite sample size & \S\ref{sec:7}--\S\ref{sec:8} \\
Place the framework in the SBI literature & \S\ref{subsec:3.7}, \S\ref{sec:10} \\
Demarcate open problems & \S\ref{sec:11} \\
\end{longtable}

Note carefully: (C1) is a \emph{frequentist} property --- it makes no
reference to a prior. The training proposal \(\rho\) on \(\Theta\) used
to draw simulations is purely a sampling distribution; in the population
limit, \emph{provided the architectural class \(\mathcal{F}\) contains
the calibration manifold} and \(\rho\) has full support on the region of
inferential interest, the CD pivot does not depend on \(\rho\). Under
misspecification or with an under-parameterized architecture, the
population minimizer of NF-MLE is the \(\rho\)-weighted KL projection
onto \(\mathcal{F}\) and \emph{does} depend on \(\rho\); see
\S\ref{subsec:11.4}.

\subsection{Scope of this document}\label{subsec:1.4}

Theory is developed in detail for the regular one-parameter case, with
explicit extensions to (i) discrete one-parameter families via auxiliary
randomization, and (ii) multivariate \(\theta\) via autoregressive
triangular flows. Misspecification is acknowledged but not addressed.

\medskip\hrule\medskip

\section{The pivot and the calibration manifold}\label{sec:2}

\subsection{Notation and standing assumptions}\label{subsec:2.1}

\begin{itemize}
\tightlist
\item
  \(\Theta \subseteq \mathbb{R}^d\) is the parameter space.
\item
  \(\mathcal{X}\) is the data space;
  \(\{P_\theta : \theta \in \Theta\}\) is the parametric statistical
  model.
\item
  \(\rho\) is the \textbf{proposal distribution} on \(\Theta\), with
  full support on the region of inferential interest and finite first
  moment. It is a sampling device for generating training data,
  \emph{not} a Bayesian prior --- it has no role at inference time.
\item
  The training distribution is
  \((\theta, X) \sim \rho \otimes P_\theta\): sample
  \(\theta \sim \rho\), then \(X \mid \theta \sim P_\theta\).
\item
  \(\varphi, \Phi\) denote the scalar standard normal density / CDF;
  \(\phi_d, \Phi_d\) their \(d\)-variate analogues (product of standard
  normals).
\item
  \(\mathrm{KL}(p \,\|\, q) := \int p \log(p/q)\,d\mu\) is the
  Kullback--Leibler divergence; \(\mathrm{KL} \ge 0\) with equality iff
  \(p = q\) a.e.
\end{itemize}

\subsection{The pivot, the implied CD, the calibration manifold}\label{subsec:2.2}

A \textbf{pivot} in the sense used here is a measurable function
\(r: \Theta \times \mathcal{X} \to \mathbb{R}^d\) satisfying:

\begin{itemize}
\tightlist
\item
  \textbf{(R1) Monotonicity in \(\theta\).} In 1D,
  \(\theta \mapsto r(\theta; X)\) is strictly increasing for every
  \(X\). In \(d > 1\), the autoregressive form
  (R1\(^\mathrm{auto}\), \S\ref{sec:6}) applies: \(r_k\) depends on
  \(\theta\) only through \(\theta_{\le k}\), and
  \(\theta_k \mapsto r_k(\theta; X)\) is strictly increasing at fixed
  \(\theta_{<k}, X\). The increasing-rather-than-decreasing direction is
  a sign convention (Lemma 4.1 admits both branches; cf.~\S\ref{subsec:4.3});
  the convention pins the orientation of \(H_r\) so that larger \(H_r\)
  corresponds to larger \(\theta\).
\item
  \textbf{(R2) Invertibility in \(X\).} For each fixed \(\theta\), the
  map \(X \mapsto r(\theta; X)\) is a \(C^1\) diffeomorphism onto its
  image, with full-rank Jacobian everywhere. The Lipschitz extension
  needed for ReLU-network arguments is handled separately in
  \S\ref{subsec:4.4} via the auxiliary condition (R3): full-rank a.e.\ is
  \emph{not} equivalent to injectivity for Lipschitz maps, so
  (R2)-by-itself is the \(C^1\) statement and the Lipschitz case requires
  the uniform lower bound \(|\partial_X r| \ge c > 0\).
\end{itemize}

The \textbf{induced CD} is \[
H_r(\theta; X) \;:=\; \Phi_d(r(\theta; X)) \;\in\; (0,1)^d,
\] and the \textbf{calibration manifold} (within a function class
\(\mathcal{F}\)) is \[
\mathcal{M}_\mathcal{F} \;:=\; \bigl\{r \in \mathcal{F} \,:\, \mathrm{Law}(r(\theta_0; X) \mid \theta_0) = \mathcal{N}(0, I_d) \text{ for every } \theta_0 \in \mathrm{supp}(\rho)\bigr\}.
\] Any element of \(\mathcal{M}_\mathcal{F}\) satisfying (R1) is a valid
CD pivot.

\textbf{Concrete example.} In the location-normal model
\(X \sim \mathcal{N}(\theta, 1)\), the function
\(r^*(\theta, X) = \theta - X\) is a pivot in this sense:
\(r^*(\theta_0, X) = \theta_0 - X = -Z \sim \mathcal{N}(0,1)\),
satisfying (R1) (\(\partial_\theta r^* = +1\)) and (R2)
(\(\partial_X r^* = -1\)). The induced CD
\(H_{r^*}(\theta; X) = \Phi(\theta - X)\) recovers the textbook one. The
framework we develop is: find the analogue of \(r^*\) when no closed
form is available, by training a neural surrogate.

\subsection{The role of (R1): why monotonicity in \(\theta\) is needed}\label{subsec:2.3}

Calibration alone gives an object with the right marginal frequentist
behavior under each \(\theta_0\), but doesn't guarantee that confidence
sets are connected or interpretable. Monotonicity in \(\theta\) ensures:

\begin{itemize}
\tightlist
\item
  \(H_r(\cdot; X)\) is a proper CDF on \(\Theta\), so the
  level-\(\alpha\) confidence interval
  \(H_r(\cdot; X)^{-1}([\alpha/2, 1-\alpha/2])\) is a connected interval
  (1D) or rectangle (multivariate via the autoregressive structure).
\item
  The ``direction'' of the CD is fixed: larger \(H_r\) corresponds to
  larger \(\theta\), matching the convention that the CD's quantile
  function is the upper endpoint of a one-sided confidence interval.
\end{itemize}

Without (R1), the calibration manifold is large: any composition of
\(r\) with a measure-preserving bijection of \(\mathcal{N}(0, I_d)\) to
itself --- sign flips in each coordinate, non-monotone measure-preserving
rearrangements in 1D, and the entire orthogonal group \(O(d)\) in
\(d > 1\) (since \(\mathcal{N}(0, I_d)\) is rotation-invariant) --- stays
in \(\mathcal{M}\). The ambiguity is infinite-dimensional in general; see
\S\ref{subsec:3.3} for the explicit characterization. (R1) in 1D
eliminates the sign-flip and non-monotone-rearrangement ambiguities. In
\(d > 1\), (R1) eliminates the coordinate-wise sign-flip ambiguity but
\emph{not} the rotational ambiguity --- the autoregressive triangular
structure of \S\ref{sec:6} is what selects the unique Knothe--Rosenblatt
representative within the rotation orbit; see \S\ref{subsec:11.1} for
the related open problem outside the autoregressive class.

\subsection{The role of (R2): why monotonicity in \(X\) is needed}\label{subsec:2.4}

Invertibility in \(X\) at fixed \(\theta\) is what makes the
change-of-variables formula valid: \[
\hat p_r(X \mid \theta) \;=\; \phi_d(r(\theta, X)) \cdot \bigl|\det \partial r / \partial X\bigr|.
\] Without this, the right-hand side is not a valid density (it misses
the sum over preimages when \(X \mapsto r\) is not bijective), and any
loss built on it is mis-specified.

When \(r\) is not injective in \(X\) at fixed \(\theta\), the
change-of-variables formula requires summing over the multiple preimages
of each value, but the surrogate \(\hat p_r\) above keeps only one
preimage's contribution. The result fails to integrate to one in \(X\)
at fixed \(\theta\); consequently \(-\log \hat p_r\) is not bounded
below by the conditional entropy of \(X \mid \theta\) (the Gibbs
inequality requires a normalized density), and for architectures that
admit unbounded \(|\det \partial r / \partial X|\) in folded regions ---
e.g.\ ReLU / UMNN networks with no built-in monotonicity --- the NF-MLE
loss has no lower bound at all. The formal version is in
\S\ref{subsec:3.5}; the empirical demonstration in \S\ref{subsec:8.4}
(where the autograd-Jacobian ablation attains loss 0.56 against the
true model's 0.88). Enforcing (R2) architecturally fixes this.

\medskip\hrule\medskip

\section*{Part II — Loss design}
\addcontentsline{toc}{section}{Part II — Loss design}

\section{The NF-MLE objective}\label{sec:3}

\subsection{What a normalizing flow is, and the statement of the loss}\label{subsec:3.1}

A \emph{normalizing flow} is a parametric family of diffeomorphisms
\(r: \mathcal{X} \to \mathbb{R}^d\) that transforms a complicated source
density \(p(X)\) into a simple base density \(\phi_d\) (we use a
standard Gaussian throughout). Conditional on \(\theta\), the implied
density on \(\mathcal{X}\) is the standard change-of-variables
expression
\(\hat p_r(X \mid \theta) = \phi_d(r(\theta, X)) \cdot |\det \partial r / \partial X|\).
\emph{NF-MLE} simply means: train \(r\) by maximum likelihood under this
implied density.

For a pivot \(r\) satisfying (R1) and (R2), define \[
\boxed{\;L(r) \;:=\; \mathbb{E}_{(\theta, X) \sim \rho \otimes P_\theta}\!\left[ \tfrac{1}{2}\|r(\theta, X)\|^2 \;-\; \log\bigl|\det \tfrac{\partial r}{\partial X}\bigr| \right].\;}
\] This is the negative log-likelihood under \(\hat p_r\) with the
irrelevant constant \(\tfrac{d}{2}\log(2\pi)\) dropped.

\subsection{Why it works: strict propriety via KL non-negativity}\label{subsec:3.2}

Write the model density implied by \(r\) as
\(\hat p_r(X \mid \theta) := \phi_d(r(\theta, X)) \cdot |\det \partial r / \partial X|\).
Direct algebra gives \[
L(r) \;=\; \mathbb{E}_\rho\bigl[\mathbb{E}_{X \mid \theta}[-\log \hat p_r(X \mid \theta)]\bigr] \;=\; \mathbb{E}_\rho\bigl[\mathrm{KL}\bigl(p(\cdot \mid \theta) \,\|\, \hat p_r(\cdot \mid \theta)\bigr)\bigr] + \mathrm{const},
\] where the constant is the conditional entropy of \(X \mid \theta\)
averaged over \(\rho\) (independent of \(r\)). Since
\(\mathrm{KL} \ge 0\) with equality iff the two densities agree a.e.,
\(L(r)\) is minimized iff
\(\hat p_r(\cdot \mid \theta) = p(\cdot \mid \theta)\) for \(\rho\)-a.e.
\(\theta\). Equivalently,
\(r(\theta_0; X) \mid \theta_0 \sim \mathcal{N}(0, I_d)\) for
\(\rho\)-a.e. \(\theta_0\). So: \[
\arg\min_{r \in \mathcal{F}} L(r) \;=\; \mathcal{M}_\mathcal{F},
\] where \(\mathcal{F}\) is the architectural class over which the
optimization ranges (assumed to satisfy (R2) so \(\hat p_r\) is a valid
density).

This is what is meant by \textbf{strict propriety}: a loss is
\emph{strictly proper} for a target \(\mathcal{M}\) if its
population-level minimum, over the class of admissible \(r\), is
attained exactly on \(\mathcal{M}\) and nowhere else. (The term comes
from the scoring-rule literature; Gneiting \& Raftery 2007 is the
standard reference.) NF-MLE is strictly proper for the calibration
manifold within the architectural class. No auxiliary independence or
sharpness penalties are needed --- the per-sample log-density already
encodes both calibration \emph{and} the implied sharpness of
\(\hat p_r\).

\subsection{The role of architectural constraints in selecting from \(\mathcal{M}\)}\label{subsec:3.3}

The calibration manifold over the \emph{unconstrained} class of pivots
is enormous. In 1D, any composition of a monotone transport
\(X \mid \theta \to \mathcal{N}(0,1)\) with a measure-preserving
bijection of \(\mathcal{N}(0,1)\) to itself remains in the manifold; the
manifold is infinite-dimensional. In multivariate, rotations of a
calibrated pivot stay calibrated since \(\mathcal{N}(0, I_d)\) is
rotation-invariant.

NF-MLE selects an element of \(\mathcal{M}\) but does not by itself
identify which one; the architectural class \(\mathcal{F}\) does. We
will see:

\begin{itemize}
\tightlist
\item
  In 1D, \emph{within the class \(\mathcal{F}_{C^1}\) of \(C^1\) pivots
  satisfying (R1) and (R2)}, \(\mathcal{M}_\mathcal{F}\) is a singleton
  (Theorem A, \S\ref{sec:4}).
\item
  In multivariate, \emph{within the autoregressive triangular class}
  satisfying the corresponding multivariate (R1), (R2),
  \(\mathcal{M}_\mathcal{F}\) reduces to a singleton: the
  Knothe--Rosenblatt rearrangement (Theorem A-d, \S\ref{subsec:6.3}).
\end{itemize}

So the slogan ``architecture picks out a unique element'' is precise:
the function-class restriction makes the population optimum unique.

\subsection{The Bayesian connection: NF-MLE is the SNL loss}\label{subsec:3.4}

The objective \(L(r)\) is \emph{identical} to the loss used by
Sequential Neural Likelihood (SNL; Papamakarios et al.~2019). The full
picture is:

\begin{longtable}[]{@{}
  >{\raggedright\arraybackslash}p{(\columnwidth - 4\tabcolsep) * \real{0.3333}}
  >{\raggedright\arraybackslash}p{(\columnwidth - 4\tabcolsep) * \real{0.3333}}
  >{\raggedright\arraybackslash}p{(\columnwidth - 4\tabcolsep) * \real{0.3333}}@{}}
\toprule\noalign{}
\begin{minipage}[b]{\linewidth}\raggedright
Component
\end{minipage} & \begin{minipage}[b]{\linewidth}\raggedright
SNL
\end{minipage} & \begin{minipage}[b]{\linewidth}\raggedright
CD-SBI
\end{minipage} \\
\midrule\noalign{}
\endhead
\bottomrule\noalign{}
\endlastfoot
Loss & NF-MLE on \(X \mid \theta\) & Same \\
Architecture & Generic normalizing flow & Flow with
monotone-in-\(\theta\) + monotone-in-\(X\) structure \\
Output & Approximate likelihood \(\hat p(X \mid \theta)\) & CD pivot
\(r(\theta; X)\) \\
Inference & MCMC on \(\hat p(X_\text{obs} \mid \theta)\,\pi(\theta)\) &
Direct inversion:
\(\{\theta : \|r(\theta; X_\text{obs})\|^2 \le \chi^2_{d,\alpha}\}\) \\
Coverage guarantee & None (Bayesian) & Pointwise frequentist \\
\end{longtable}

In the population limit on a sufficiently expressive class, both SNL and
CD-SBI recover the same conditional density \(p(X \mid \theta)\) at the
NF-MLE optimum. CD-SBI's contribution is the architectural prescription
that \emph{(a)} exposes a monotone-in-\(\theta\) pivot directly, and
\emph{(b)} makes pivot inversion well-defined. The architectural
constraints do not change what is learned in regular cases (where the
true pivot already satisfies monotonicity); they restrict the
parameterization to make the CD output well-typed.

This is good news: NF-MLE inherits all of SNL's optimization stability,
scalability, and tooling.

\subsection{Why architectural monotonicity in \(X\) is essential}\label{subsec:3.5}

If \(r\) is not bijective in \(X\) at fixed \(\theta\), the
change-of-variables formula does not apply, and \(\hat p_r\) is not a
valid density. NF-MLE then becomes ill-posed. In practice --- using
autograd to compute the Jacobian during training --- the model can find
regions of \(X\)-space where \(r\) folds and the local
\(|\partial r / \partial X|\) becomes very large; this gives an
artificially low loss while \(\hat p_r\) is not actually matching \(p\).
Because the folded \(\hat p_r\) is not normalized in \(X\) at fixed
\(\theta\), the loss \(-\mathbb{E}_X[\log \hat p_r]\) has no lower bound
at the conditional entropy, so it can fall \emph{below} the loss
attained by the true \(r^*\). The empirical demonstration is in \S\ref{subsec:8.4}.

The clean fix is to enforce strict monotonicity in \(X\) architecturally
--- either via \emph{Unconstrained Monotonic Neural Networks} (UMNNs;
Wehenkel \& Louppe 2019), which represent a scalar monotone function as
the integral of a positive function from a free neural net, or via
autoregressive triangular structure. In both cases the Jacobian is
closed-form and positive by construction.

\medskip\hrule\medskip

\setcounter{subsection}{6}
\subsection{The design space of strictly proper objectives}\label{subsec:3.7}

The structure of the problem --- find a function \(r\) whose minimizers,
over a suitable class, are exactly the calibration manifold
\(\mathcal{M}\) --- admits multiple families of candidate losses. This
subsection situates NF-MLE within that design space and contrasts with
the choices made by other SBI methods.

\emph{Brief glossary for what follows.} The \textbf{Cramér--von Mises
(CvM)} distance is \(\int (F_n(z) - F_0(z))^2\, dF_0(z)\), a strictly
proper two-sample distance between an empirical CDF \(F_n\) and a
reference CDF \(F_0\). The \textbf{energy distance} is a related
\(L^2\)-distance between CDFs in higher dimensions. \textbf{HSIC}
(Hilbert--Schmidt Independence Criterion; Gretton et al.~2005) is a
kernel-based test statistic that is zero iff two random variables are
independent. The \textbf{continuous ranked probability score (CRPS)} is
\(\int(F(z) - \mathbf{1}\{y \le z\})^2\,dz\), a strictly proper scoring
rule for predictive distributions; the \textbf{pinball loss} is its
quantile analogue.

\textbf{Class 1: Two-sample distance on the marginal pushforward.} Apply
any strictly proper distance --- CvM, energy distance, KL --- between
\(\mathrm{Law}(\Phi(r(\theta, X)))\) and \(U(0,1)^d\) on the joint
training distribution. Strictly proper for \emph{marginal} calibration.
\textbf{Insufficient on its own}: marginal \(\Phi(r) \sim U(0,1)\) does
not imply conditional \(\Phi(r) \mid \theta \sim U(0,1)\) --- the
Hermans trust-crisis pattern lives precisely here, with marginal PIT
passing while conditional fails.

\textbf{Class 2: Marginal + independence.} Augment Class 1 with an
independence term \(\mathrm{HSIC}(\theta, \Phi(r))\) or similar.
Marginal calibration plus \(\theta \perp \Phi(r)\) implies conditional
calibration. Strictly proper for \(\mathcal{M}\), but empirically
fragile in three ways: (i) the HSIC term has a weak gradient signal far
from the optimum compared to a per-sample density target; (ii) the
marginal-calibration term and the independence term must be balanced via
hyperparameters, with no principled choice; (iii) the marginal and
conditional terms can decouple in finite samples, reproducing the
Class-1 pathology.

\textbf{Class 3: Stratified / conditional matching.} Discretize
\(\Theta\) into bins; within each bin, apply a two-sample distance
between \(\Phi(r(\theta, X)) \mid \theta \in \mathrm{bin}\) and
\(U(0,1)^d\). Strictly proper for each conditional. \textbf{Doesn't
scale}: requires \(K^d\) bins in \(d\)-dimensional \(\Theta\), defeating
the entire motivation for amortized inference.

\textbf{Class 4: Scoring rules on standardized residuals.} Apply a
strictly proper scoring rule --- CRPS, pinball --- to \(r(\theta, X)\)
against the reference \(\mathcal{N}(0, I_d)\). \textbf{Directional
propriety error.} Recall the standard form: a scoring rule \(S(F, y)\)
is strictly proper if for any \(G^*\), the expectation
\(\mathbb{E}_{Y \sim G^*}[S(F, Y)]\) over the \emph{sample} is minimized
over the \emph{forecast} \(F\) at \(F = G^*\). Here the situation is
reversed: the reference \(F = \mathcal{N}(0, I_d)\) is fixed and
\(Y = r(\theta, X)\) is the learned quantity, so we are minimizing over
the \emph{sample distribution} \(G\) with \(F\) fixed, which is not the
proper direction. Concretely, for univariate CRPS with fixed \(F\) and
\(Y \sim G\), \[
\mathbb{E}_{Y \sim G}[\mathrm{CRPS}(F, Y)] \;=\; \int F(z)^2\,dz \;+\; \int G(z)\bigl(1 - 2F(z)\bigr)\,dz,
\] since \(\mathbb{E}[\mathbf{1}\{Y \le z\}] = G(z)\) and the cross-term
yields the displayed form. The first integral does not depend on \(G\).
For \(F = \Phi\), the coefficient \(1 - 2\Phi(z)\) is positive on
\(z < 0\) and negative on \(z > 0\), so the pointwise-optimal \(G\)
satisfies \(G(z) = 0\) for \(z < 0\) and \(G(z) = 1\) for \(z > 0\) ---
i.e., \(G = \delta_0\). The CRPS-against-\(\mathcal{N}(0,1)\) loss is
\emph{minimized by collapsing \(r(\theta, X)\) to a constant at the
median}, the opposite of calibration. (This is why LF2I uses
pinball-style losses in the \emph{opposite} direction --- fixing the
test statistic and learning critical values --- and is not a
counterexample.)

\textbf{Class 5: Per-sample conditional log-density (NF-MLE).} Apply KL
on the conditional density:
\(L(r) = \mathbb{E}_{(\theta, X)}\bigl[-\log \hat p_r(X \mid \theta)\bigr]\)
where \(\hat p_r\) is the change-of-variables density implied by \(r\).
Strictly proper for \(\mathcal{M}_\mathcal{F}\) by \S\ref{subsec:3.2}. Scales
pointwise to any dimension. Identical to the SNL loss. \textbf{The only
subtlety is the architectural monotonicity-in-\(X\) requirement} (\S\ref{subsec:3.5})
--- without it, \(\hat p_r\) is not a valid density and the propriety
argument breaks.

\subsubsection*{Position in the SBI literature}

The five classes intersect existing SBI methods as follows. NPE/SNPE
(Papamakarios \& Murray 2016; Greenberg et al.~2019) minimize per-sample
log-density of \(\theta \mid X\) --- a Class-5 loss on the
\emph{posterior}, with no built-in coverage. NLE/SNL (Papamakarios et
al.~2019) minimize per-sample log-density of \(X \mid \theta\) ---
Class-5 on the \emph{likelihood}, which is the same loss CD-SBI uses;
the difference is the architectural prescription (monotone in \(\theta\)
+ monotone in \(X\)). NRE (Hermans et al.~2020) targets the ratio via
binary classification --- a Class-5-flavored loss on a different object.
Balanced NRE (Delaunoy et al.~2022) and Calibrated NPE (Falkiewicz et
al.~2023) add regularization terms to address conservativeness post hoc;
CD-SBI achieves calibration by construction instead. LF2I (Dalmasso et
al.~2024), WALDO (Masserano et al.~2023), and Box CD (Bortolato \&
Ventura 2025) operate in a structurally different mode --- they learn
test statistics (or hyper-rectangular acceptance regions) and invert via
Neyman construction; the calibration target is the CDF of the test
statistic under each \(\theta\), not the CDF of the parameter under
data. CD-SBI's contribution within this taxonomy is the architectural
recipe that converts Class-5/SNL into a procedure whose output is a
frequentist CD by construction.

\medskip\hrule\medskip

\section*{Part III — Theory in one parameter}
\addcontentsline{toc}{section}{Part III — Theory in one parameter}

\section{Uniqueness in the location-normal model}\label{sec:4}

We work in 1D (\(d = 1\)) and develop the theory for
\(X \sim \mathcal{N}(\theta, 1)\), then lift to general regular
exponential families in \S\ref{sec:5}.

\subsection{Statement}\label{subsec:4.1}

\textbf{Theorem A (Uniqueness, \(C^1\)).} \emph{Let
\(X \sim \mathcal{N}(\theta, 1)\) on \(\Theta = \mathbb{R}\). Within the
class} \[
\mathcal{F}_{C^1} \;:=\; \bigl\{ r : \mathbb{R}^2 \to \mathbb{R} \text{ satisfying (R1), (R2), with } X \mapsto r(\theta_0; X) \in C^1(\mathbb{R}) \text{ for each } \theta_0 \bigr\},
\] \emph{the calibration manifold is a singleton:} \[
\mathcal{M} \cap \mathcal{F}_{C^1} \;=\; \{r^*\}, \qquad r^*(\theta; X) = \theta - X.
\] \emph{Equivalently, \(L(r)\) over \(\mathcal{F}_{C^1}\) has a unique
minimizer, \(r = r^*\), with \(H_{r^*}(\theta; X) = \Phi(\theta - X)\)
--- the Schweder--Hjort UMP-unbiased CD.}

The proof rests on three lemmas.

\subsection{The three lemmas}\label{subsec:4.2}

\textbf{Lemma 4.1 (Monotone transport).} \emph{Let
\(\theta_0 \in \mathbb{R}\) and \(\psi: \mathbb{R} \to \mathbb{R}\)
strictly monotone such that \(\psi\) pushes \(\mathcal{N}(\theta_0, 1)\)
forward to \(\mathcal{N}(0,1)\). Then \(\psi(X) = X - \theta_0\) (if
strictly increasing) or \(\psi(X) = \theta_0 - X\) (if strictly
decreasing).}

\emph{Proof.} Standard 1D monotone rearrangement: between two atomless
probability measures on \(\mathbb{R}\), the unique monotone-increasing
measure-preserving map is the composition of CDF and inverse CDF. For
strictly increasing \(\psi\): \(\Phi(\psi(X)) = \Phi(X - \theta_0)\)
a.s. by uniqueness, so \(\psi(X) = X - \theta_0\). For strictly
decreasing \(\psi\),
\(\Phi(\psi(X)) = 1 - \Phi(X - \theta_0) = \Phi(\theta_0 - X)\), so
\(\psi(X) = \theta_0 - X\). \(\blacksquare\)

\textbf{Lemma 4.2 (Level-set rigidity, \(C^1\)).} \emph{Let
\(\theta_0 \in \mathbb{R}\) and \(\psi \in C^1(\mathbb{R})\) pushing
\(\mathcal{N}(\theta_0, 1)\) to \(\mathcal{N}(0,1)\). Then \(\psi'\) has
constant sign on \(\mathbb{R}\); i.e., \(\psi\) is strictly monotone.}

\emph{Proof.} Suppose for contradiction \(\psi'\) changes sign on
\(\mathbb{R}\). By continuity, there is some \(X_*\) at which \(\psi'\)
crosses zero, and \(\psi\) attains a strict local extremum at \(X_*\)
(WLOG a local maximum with \(v_* := \psi(X_*)\); the local-minimum case
is symmetric). For \(v\) just below \(v_*\), the level set
\(\{X : \psi(X) = v\}\) contains at least two points \(X_1(v) \to X_*^-\)
and \(X_2(v) \to X_*^+\), with \(\psi\) strictly monotone on each of
\((X_1, X_*)\) and \((X_*, X_2)\).

The pushforward density of \(\mathcal{N}(\theta_0, 1)\) under \(\psi\)
at any value \(v\) with finite preimage is
\[
g_\psi(v) \;=\; \sum_{i\,:\,\psi(X_i)=v} \frac{\varphi(X_i - \theta_0)}{|\psi'(X_i)|},
\]
so summing over the two preimages above,
\[
g_\psi(v) \;\geq\; \frac{\varphi(X_1(v) - \theta_0)}{|\psi'(X_1(v))|} + \frac{\varphi(X_2(v) - \theta_0)}{|\psi'(X_2(v))|}.
\]
Let \(2m \ge 2\) be the order of the leading nonzero derivative of
\(\psi\) at \(X_*\), so \(\psi(X) - v_* \sim -c\,(X - X_*)^{2m}\) near
\(X_*\) with \(c > 0\). Then
\(|\psi'(X_i(v))| \sim |v - v_*|^{(2m-1)/(2m)}\), while
\(\varphi(X_i(v) - \theta_0) \to \varphi(X_* - \theta_0) > 0\). Hence
\(g_\psi(v)\) diverges at least as fast as \(|v - v_*|^{-1/2}\) (and
faster for higher-order critical points). But the target density
\(\varphi(v_*)\) is finite, a contradiction. \(\blacksquare\)

\emph{Remark (general source).} The proof uses only two facts about
\(\mathcal{N}(\theta_0, 1)\) and \(\mathcal{N}(0,1)\): the source has
positive density at any critical point of \(\psi\), and the target has
finite density at all attained values. The lemma therefore applies
verbatim when \(\mathcal{N}(\theta_0, 1)\) is replaced by any continuous
distribution on \(\mathbb{R}\) with positive density on its support and
\(\mathcal{N}(0,1)\) by any continuous distribution with finite density.
This generalization is invoked silently in Theorem~C (\S\ref{subsec:5.2}, with the
pushforward of \(T(X) \mid \theta_0\) as source) and in the inductive
step of Theorem~A-d (\S\ref{subsec:6.3}, with the conditional law
\(F_k^{(\theta)}(\cdot \mid X_{<k})\) as source).

\textbf{Lemma 4.3 (NF-MLE is strictly proper, conditional).} \emph{Under
the NF-MLE objective \(L(r)\), any minimizer \(r \in \mathcal{F}_{C^1}\)
satisfies \(r(\theta_0; X) \mid \theta_0 \sim \mathcal{N}(0,1)\) for
\(\rho\)-a.e. \(\theta_0\).}

\emph{Proof.} From \S\ref{subsec:3.2}, \(L(r)\) is the expected KL between
\(p(\cdot \mid \theta)\) and \(\hat p_r(\cdot \mid \theta)\), up to a
constant. KL is zero iff
\(\hat p_r(\cdot \mid \theta) = p(\cdot \mid \theta)\) a.e. on
\(\mathrm{supp}\,\rho\). Under (R2), \(\hat p_r(\cdot \mid \theta)\)
pushes \(X \mid \theta\) to \(\mathcal{N}(0,1)\) iff
\(r(\theta_0; X) \mid \theta_0 \sim \mathcal{N}(0,1)\). \(\blacksquare\)

\emph{Remark (Lipschitz / autoregressive variants).} The proof depends
only on (R2) and the KL identity of \S\ref{subsec:3.2}; it therefore
extends without change to any pivot class for which the change-of-variables
formula yields a valid normalized density --- in particular to the
Lipschitz form of (R2) used in Theorem A* (\S\ref{subsec:4.4}) and to the
autoregressive variant (R2\(^{\mathrm{auto}}\)) used in \S\ref{sec:6}.

\subsection{Proof of Theorem A}\label{subsec:4.3}

\emph{Existence.} \(r^*(\theta; X) = \theta - X\) is \(C^\infty\). Under
\(X = \theta_0 + Z\) with \(Z \sim \mathcal{N}(0,1)\),
\(r^*(\theta_0; X) = -Z \sim \mathcal{N}(0,1)\), so
\(r^* \in \mathcal{M}\). And \(\partial_\theta r^* = 1 > 0\) (R1),
\(\partial_X r^* = -1 \neq 0\) (R2).

\emph{Uniqueness.} Let \(r \in \mathcal{M} \cap \mathcal{F}_{C^1}\). For
each \(\theta_0 \in \mathbb{R}\), define
\(\psi_{\theta_0}: X \mapsto r(\theta_0; X)\). By Lemma 4.3,
\(\psi_{\theta_0}\) pushes \(\mathcal{N}(\theta_0, 1)\) to
\(\mathcal{N}(0,1)\). By Lemma 4.2 (\(\psi_{\theta_0}\) is \(C^1\) by
hypothesis), \(\psi_{\theta_0}\) is strictly monotone. By Lemma 4.1,
\(\psi_{\theta_0}(X) \in \{X - \theta_0,\, \theta_0 - X\}\).

\emph{Selecting the direction.} A priori, (R1) does not rule out a
measurable mixture in which \(\psi_{\theta_0}(X) = X - \theta_0\) on
some set \(A \subseteq \Theta\) and \(\psi_{\theta_0}(X) = \theta_0 - X\)
on \(A^c\). But for any fixed \(X\), such a mixture would make
\(\theta \mapsto r(\theta, X) = \pm(\theta - X) + (\mathrm{const})\)
decreasing on \(A\) and increasing on \(A^c\), violating the global
strict monotonicity required by (R1). Hence one branch holds
\(\rho\)-a.e. The increasing-branch case \(\psi_{\theta_0}(X) = X - \theta_0\)
gives \(r(\theta, X) = X - \theta\) with \(\partial_\theta r = -1\),
contradicting (R1). The other branch
\(\psi_{\theta_0}(X) = \theta_0 - X\) gives \(r(\theta, X) = \theta - X\)
with \(\partial_\theta r = +1\), satisfying (R1). The full-support
assumption on \(\rho\) (\S\ref{subsec:2.1}) lifts this from \(\rho\)-a.e.
\(\theta_0\) to all of \(\mathbb{R}\). \(\blacksquare\)

\subsection{Strengthening to Lipschitz: Theorem A*}\label{subsec:4.4}

Lemma 4.2 uses \(C^1\) regularity essentially: the contradiction comes
from a \(|v - v_*|^{-1/2}\) (or faster) divergence of the pushforward
density near a critical value, which fails for Lipschitz V-shapes. For
ReLU networks (Lipschitz, piecewise affine), we need an alternative.

The clean approach is to require strict \(X\)-monotonicity as an extra
condition:

\begin{itemize}
\tightlist
\item
  \textbf{(R3) {[}strict \(X\)-monotonicity, a.e.{]}:} \(\partial_X r\)
  has constant sign and is bounded away from zero on the
  differentiability set: \(|\partial_X r(\theta; X)| \geq c > 0\) a.e.
\end{itemize}

\textbf{Theorem A* (Uniqueness, Lipschitz).} \emph{Suppose
\(r: \mathbb{R}^2 \to \mathbb{R}\) is Lipschitz and satisfies (R1),
(R2), (R3). If \(r \in \mathcal{M}\) then \(r(\theta; X) = \theta - X\)
a.e.}

\emph{Proof.} By Rademacher's theorem (a Lipschitz function on
\(\mathbb{R}^n\) is differentiable Lebesgue-a.e.), \(r\) is
differentiable a.e. By (R3), \(\psi_{\theta_0}(X) := r(\theta_0; X)\) is
strictly monotone in \(X\). By Lemma 4.3, \(\psi_{\theta_0}\) pushes
\(\mathcal{N}(\theta_0, 1)\) to \(\mathcal{N}(0,1)\). By Lemma 4.1 (1D
monotone rearrangement, valid for any pair of atomless probability
measures on \(\mathbb{R}\)),
\(\psi_{\theta_0}(X) \in \{X - \theta_0,\, \theta_0 - X\}\) a.e. Sign
determination as in Theorem A. \(\blacksquare\)

When is (R3) automatic? Whenever the architecture enforces it. For
\(C^1\) networks (tanh, GELU, softplus activations), Lemma 4.2 derives
it from \(C^1\) + calibration. For UMNN architectures with \(X\) as a
monotone input, by construction. For triangular flows of the kind
described in \S\ref{sec:6}, by construction. For vanilla ReLU networks with no
architectural enforcement, (R3) may fail, and architectural enforcement
is the recommended fix.

\subsection{The Schweder--Hjort bridge}\label{subsec:4.5}

A \textbf{UMP-unbiased CD} is, in the classical sense, the CD whose
associated one-sided test family at each level is uniformly most
powerful within the class of unbiased tests; this is the analogue of the
UMPU confidence interval in classical testing theory. For
\(X \sim \mathcal{N}(\theta, 1)\), Schweder \& Hjort (2016, ch.~5) give
the UMP-unbiased CD as
\(C^*_\mathrm{SH}(\theta; X) = \Phi(\theta - X)\). Our optimal pivot
\(r^*(\theta; X) = \theta - X\) induces
\(H_{r^*}(\theta; X) = \Phi(\theta - X) = C^*_\mathrm{SH}\). So:

\textbf{Proposition 4.5.} \emph{In the location-normal model, the unique
element of \(\mathcal{M} \cap \mathcal{F}_{C^1}\) produced by CD-SBI is
identically the Schweder--Hjort UMP-unbiased CD.}

This is what we mean by ``the framework reduces to the classical theory
in regular cases.''

\medskip\hrule\medskip

\section{Lift to regular one-parameter exponential families}\label{sec:5}

\subsection{Setup and canonical pivot}\label{subsec:5.1}

\emph{Background, brief.} A \emph{one-parameter exponential family} is a
parametric model whose density takes the canonical form
\(p(X \mid \theta) = h(X) \exp(\theta \cdot T(X) - A(\theta))\), with
\(T(X)\) a real-valued \textbf{sufficient statistic} for \(\theta\) ---
meaning \(X\) enters the likelihood only through \(T(X)\). The family
has \textbf{monotone likelihood ratio} (MLR) in \(T\) if for any
\(\theta_1 < \theta_2\), \(p(X \mid \theta_2)/p(X \mid \theta_1)\) is a
monotone function of \(T(X)\). This is automatic in the canonical form
when \(T\) is real-valued. MLR is the hypothesis that powers most
classical optimal-testing results.

Now let \(\{P_\theta : \theta \in \mathbb{R}\}\) be a regular
one-parameter exponential family with sufficient statistic
\(T : \mathcal{X} \to \mathbb{R}\). Assume \(T(X)\) has a continuous
positive density under each \(P_\theta\), and assume increasing MLR.
(The decreasing-MLR case differs only in sign convention.)

The \textbf{canonical CD pivot} is \[
r^*(\theta; X) \;:=\; \Phi^{-1}\bigl(1 - F_\theta(T(X))\bigr),
\] where \(F_\theta\) is the CDF of \(T(X)\) under \(P_\theta\). (For
decreasing MLR, drop the \(1 -\).)

\emph{Verification.} Under \(X \sim P_{\theta_0}\),
\(F_{\theta_0}(T(X)) \sim U(0,1)\) by the probability integral
transform, so \(r^*(\theta_0; X) \sim \mathcal{N}(0,1)\). (R1) holds by
MLR (the CDF \(F_\theta\) is monotone-decreasing in \(\theta\) at any
fixed \(t\) under increasing MLR). (R2) holds since \(r^*\) is strictly
monotone in \(T(X)\) and we assume the density of \(T(X)\) is positive
on its support.

\subsection{Theorem C: uniqueness}\label{subsec:5.2}

\textbf{Theorem C (Lift to exponential families).} \emph{Under the setup
of \S\ref{subsec:5.1}, within the class
\(\mathcal{F}^T_{C^1} := \{r : t \mapsto r(\theta_0; T^{-1}(t)) \in C^1 \text{ for each } \theta_0\}\),
we have \(\mathcal{M} \cap \mathcal{F}^T_{C^1} = \{r^*\}\).}

\emph{Proof sketch.} Apply Theorem A in \(t = T(X)\) coordinates. The
pushforward density of \(T(X) \mid \theta_0\) is continuous and positive
on its support. The same three-lemma argument applies, with the source
density playing the role of the Gaussian. Direction is fixed by MLR.
\(\blacksquare\)

\subsection{Worked examples}\label{subsec:5.3}

\emph{An \textbf{ancillary statistic} is one whose distribution does not
depend on the parameter; conditioning on an ancillary loses no
information and is the standard tool for reducing nuisance dimensions in
classical inference (Cox 1958).}

\textbf{Location-normal with unknown variance.}
\(X_i \sim \mathcal{N}(\theta, \sigma^2)\) iid, both unknown.
Conditioning on the ancillary \(S^2 := \sum_i (X_i - \bar X)^2\), the
conditional model is a one-parameter location problem with pivot \[
r^*(\theta; \bar X, S) \;=\; \Phi^{-1}\bigl( F_{t_{n-1}}\bigl(\sqrt{n}\,(\theta - \bar X)/S\bigr) \bigr).
\] This is the Student-\(t\) CD.

\textbf{Exponential rate.} \(X_i \sim \mathrm{Exp}(\theta)\) iid,
\(\theta > 0\) rate. Sufficient statistic
\(T = \sum_i X_i \sim \mathrm{Gamma}(n, \theta)\), with
\(2\theta T \sim \chi^2_{2n}\). MLR is decreasing in \(\theta\) (larger
\(\theta\) $\Rightarrow$ smaller mean of \(T\)), so \[
r^*(\theta; T) \;:=\; \Phi^{-1}\bigl( F_{\chi^2_{2n}}(2\theta T)\bigr).
\] This is the test case for whether the framework handles non-additive
pivots (it does, with the right architecture; see \S\ref{subsec:8.4}).

\medskip\hrule\medskip

\setcounter{subsection}{6}
\subsection{Discrete extension via randomization}\label{subsec:5.7}

For discrete sufficient statistics (T(X) countable), no function of
\(T(X)\) alone can have a continuous distribution, so the strict
calibration target \(\mathcal{N}(0,1)\) is unattainable. The standard
fix (Lancaster 1961) is auxiliary randomization: smear each atom across
an interval of an auxiliary uniform variable, so that the combined
object is continuous.

\subsubsection{Randomized PIT}\label{subsubsec:5.7.1}

Let \(X \sim P_\theta\) discrete with sufficient statistic \(T(X)\) and
MLR in \(\theta\). Define \(F_\theta(t) = P_\theta(T \leq t)\),
\(F_\theta(t^-) = P_\theta(T < t)\), \(p_\theta(t) = P_\theta(T = t)\).
Introduce \(U \sim U(0,1)\) independent of \(X\). The \textbf{randomized
PIT} is \[
V_\theta(X, U) \;:=\; F_\theta(T(X)^-) + U \cdot p_\theta(T(X)).
\] \emph{Fact.} Under \(X \sim P_{\theta_0}\), \(U \sim U(0,1)\)
independent: \(V_{\theta_0}(X, U) \sim U(0,1)\) exactly.

The CD-SBI framework lifts to this setting by enlarging the input space:
the pivot becomes
\(r: \Theta \times \mathcal{X} \times [0,1] \to \mathbb{R}\), training
data are triples \((\theta_i, X_i, U_i)\), and the canonical pivot is \[
r^*_\mathrm{rand}(\theta; X, U) \;:=\; \Phi^{-1}(1 - V_\theta(X, U)).
\]

\subsubsection{Theorem C* (Discrete uniqueness)}\label{subsubsec:5.7.2}

Uniqueness in the randomized class requires four regularity conditions:
(R1) monotone in \(\theta\), (R2) invertible in \((X, U)\), (R3\(_U\))
strict monotone in \(U\) (distinct from (R3) of \S\ref{subsec:4.4}: that
condition was a uniform lower bound \(|\partial_X r| \ge c > 0\) used in
the Lipschitz case; (R3\(_U\)) here is the strict-monotone-in-\(U\)
condition needed inside each \(T\)-atom in the discrete setting), plus a
new condition:

\begin{itemize}
\tightlist
\item
  \textbf{(R4) Order-preservation in \(T(X)\):} for fixed
  \((\theta, U)\), \(r(\theta; X_1, U) \le r(\theta; X_2, U)\) whenever
  \(T(X_1) < T(X_2)\), with strict inequality somewhere.
\end{itemize}

\textbf{Theorem C* (Discrete).} \emph{Under (R1)--(R4), the unique pivot
in the randomized calibration manifold is \(r^*_\mathrm{rand}\).}

The proof is a ``lex-order monotone rearrangement'': the source measure
\(P_\theta \otimes U(0,1)\) on \(\mathrm{supp}(T) \times [0,1]\) has, in
lexicographic order \((T, U)\), a unique monotone measure-preserving map
to \(U(0,1)\), namely \(V^*_\theta\). (R4) is what fixes the lex order
across \(T\)-values; without it, ``interleaving'' counterexamples exist.

\textbf{Explicit counterexample.} Take
\(X \sim \mathrm{Bin}(1, \theta)\), so \(T(X) = X \in \{0,1\}\),
\(p_\theta(0) = 1-\theta\), \(p_\theta(1) = \theta\). The canonical
\(V_\theta\) assigns \(X = 0\) to the sub-interval \([0, 1-\theta]\) of
\((0,1)\) and \(X = 1\) to \([1-\theta, 1]\): \[
V_\theta^*(0, U) = U(1-\theta), \qquad V_\theta^*(1, U) = (1-\theta) + U\theta.
\] Both are monotone in \(U\), both have output in \((0,1)\), and the
mixture \((1-\theta)\,V_\theta^*(0, U) + \theta\,V_\theta^*(1, U)\)
(sampling \(X\) then \(U\)) is uniform.

Now define an alternative \(V'_\theta\) that swaps the lex order ---
assigning \(X = 0\) to \emph{two} disjoint sub-intervals (one in the
lower half, one in the upper half) and \(X = 1\) to their complement: \[
V'_\theta(0, U) = \begin{cases} \tfrac{1-\theta}{2} \cdot 2U & \text{if } U \in [0, \tfrac{1}{2}), \\ \tfrac{1}{2} + \tfrac{1-\theta}{2} \cdot 2(U - \tfrac{1}{2}) & \text{if } U \in [\tfrac{1}{2}, 1], \end{cases}
\] with \(V'_\theta(1, U)\) defined so that the joint mixture remains
uniform on \((0,1)\) (by filling the complementary sub-intervals). Both
\(V_\theta^*\) and \(V'_\theta\) satisfy (R1)--(R3\(_U\)) and both yield
calibrated pivots via \(r = \Phi^{-1}(1 - V)\). They differ in the
lex-ordering of \(T\)-values relative to \(U\): \(V_\theta^*\) places
all of \(X=0\) below all of \(X=1\), while \(V'_\theta\) interleaves
them. (R4) selects \(V_\theta^*\).

\textbf{Why (R4) is automatic in the continuous case.} In the
continuous case, Theorem C (\S\ref{subsec:5.2}) --- which lifts Theorem A
to \(t = T(X)\) coordinates via the general-source form of Lemma~4.2
(remark, \S\ref{subsec:4.2}) --- already concludes that the canonical
pivot \(\tilde r(\theta; t)\) is monotone-increasing in \(t\) under
increasing MLR. This is exactly (R4) at the \(t\)-level. (R4) becomes
a non-redundant assumption only in the discrete setting, where monotone
rearrangement of an atomic source onto a continuous target leaves a
permutation freedom across the \(T\)-atoms that the continuous-case
rigidity (Lemma~4.2 in \(C^1\), or strict-monotonicity-in-\(X\) under
(R3) in the Lipschitz case) eliminates automatically.

\subsubsection{Mid-p alternative}\label{subsubsec:5.7.3}

A deterministic alternative is the ``mid-p'' CD, fixing \(U = 1/2\): \[
H^\mathrm{mid}_\theta(X) \;:=\; 1 - V_\theta(X, 1/2).
\] This is approximately calibrated, with total-variation error from
exact \(U(0,1)\) bounded by
\(\tfrac{1}{2}\sup_{\theta,t} p_\theta(t) = O(1/\sqrt{n})\) for
\(\mathrm{Bin}(n,\theta)\) at moderate \(\theta\). For applications
where ease of implementation outweighs the small calibration error, the
mid-p variant trains \(r(\theta; X)\) directly with \(U\) fixed at
\(1/2\), dropping the auxiliary input. See Hwang \& Yang (2001) for the
optimality theory of mid-p in the closely related contingency-table
setting.

\medskip\hrule\medskip

\section*{Part IV — Multivariate generalization}
\addcontentsline{toc}{section}{Part IV — Multivariate generalization}

\section{The triangular-autoregressive flow construction}\label{sec:6}

For \(d > 1\), the architectural recipe is \emph{autoregressive
triangular flows}. An autoregressive flow is one in which the \(k\)-th
output depends only on inputs with index \(\le k\), so the Jacobian
matrix is lower-triangular and its determinant is the product of
diagonal entries. This is the multivariate generalization of (R1) + (R2)
into a structure that scales linearly in \(d\) and admits closed-form
Jacobians. Standard references on this architectural pattern include
Papamakarios et al.~(2021) and Durkan et al.~(2019).

\subsection{The architectural class}\label{subsec:6.1}

Let
\(\theta = (\theta_1, \ldots, \theta_d) \in \Theta \subseteq \mathbb{R}^d\)
and \(X \in \mathcal{X} \subseteq \mathbb{R}^d\) (or higher-dim \(X\)
with \(d\)-dim \(r\)). We work in the class of pivots of the form
\(r = (r_1, \ldots, r_d)\) where each component obeys:

\begin{itemize}
\tightlist
\item
  \textbf{(R1\(^\mathrm{auto}\)) Autoregressive monotonicity in
  \(\theta\):} \(r_k\) depends on \(\theta\) only through
  \(\theta_{\le k}\), and \(\theta_k \mapsto r_k\) is strictly
  increasing at fixed \(\theta_{<k}, X\).
\item
  \textbf{(R2\(^\mathrm{auto}\)) Autoregressive monotonicity in \(X\):}
  \(r_k\) depends on \(X\) only through \(X_{\le k}\), and
  \(X_k \mapsto r_k\) is strictly monotone at fixed \(\theta, X_{<k}\).
\end{itemize}

(R1\(^\mathrm{auto}\)) and (R2\(^\mathrm{auto}\)) are the autoregressive
specializations of (R1) and (R2). (R1\(^\mathrm{auto}\)) implies the
coordinate-wise version of (R1) by inspection; (R2\(^\mathrm{auto}\))
implies (R2) (a lower-triangular Jacobian with strictly monotone
diagonal entries has full rank, so \(r\) is a local diffeomorphism in
\(X\) at each \(\theta\)). The strict-propriety argument of
\S\ref{subsec:3.2} therefore applies verbatim in the autoregressive
class, with no separate derivation needed.

This class admits at least two practical parameterizations:

\begin{enumerate}
\def\labelenumi{\arabic{enumi}.}
\tightlist
\item
  \textbf{Additive form (Gaussian-like models, \S\ref{subsec:8.1}--8.3):}
  \[r_k(\theta, X) = a_k(\theta_k;\, \theta_{<k}, X_{<k}) \;-\; b_k(X_k;\, \theta_{<k}, X_{<k}),\]
  with \(a_k\) a UMNN monotone in \(\theta_k\) and \(b_k\) a UMNN
  monotone in \(X_k\), both conditioned on the autoregressive context.
\item
  \textbf{Doubly-monotone UMNN form (non-additive interactions, \S\ref{subsec:8.4}):}
  \[r_k(\theta, X) = b_k(X_k;\, X_{<k}) \;+\; \int_{\theta_{k,\mathrm{ref}}}^{\theta_k} \mathrm{softplus}\bigl(\alpha_k(t;\, \theta_{<k}, X_{<k}) + \beta_k(X_k;\, X_{<k})\bigr)\, dt,\]
  which is monotone in both \(\theta_k\) (by positivity of the
  integrand) and \(X_k\) (by appropriate UMNN sub-structure on \(b_k\)
  and \(\beta_k\)).
\end{enumerate}

Both forms satisfy (R1\(^\mathrm{auto}\)) and (R2\(^\mathrm{auto}\));
the Knothe--Rosenblatt uniqueness theorem below applies to both.

\subsection{Jacobian structure and computational properties}\label{subsec:6.2}

In either parameterization, both Jacobians are lower-triangular by
construction: \[
\bigl[\partial_\theta r\bigr]_{k\ell} = \begin{cases} \partial_{\theta_k} r_k > 0 & k = \ell \\ \text{(free)} & \ell < k \\ 0 & \ell > k \end{cases}
\qquad
\bigl[\partial_X r\bigr]_{k\ell} = \begin{cases} \mathrm{sgn} \cdot \partial_{X_k} r_k & k = \ell \\ \text{(free)} & \ell < k \\ 0 & \ell > k \end{cases}
\] (with a consistent sign convention on \(\partial_X r\) across
coordinates). The determinant is the product of diagonals, and \[
\log\bigl|\det \partial_X r\bigr| \;=\; \sum_{k=1}^{d} \log \bigl|\partial_{X_k} r_k\bigr|,
\] all closed-form from the UMNN positive integrand. No
autograd-on-autograd is needed; this is essential for stable training in
\(d > 1\).

\textbf{Connected CD level sets.} \(r\) is monotone-increasing in
\(\theta_k\) at fixed \(\theta_{<k}\), so \(H_r(\cdot; X)\) has
connected level sets in the natural autoregressive sense. Confidence
sets via the chi-square inversion are well-defined.

\textbf{Linear scaling in \(d\).} All operations cost \(O(d)\) per
sample, vs.~\(O(d^3)\) for a general dense Jacobian.

\subsection{Multivariate uniqueness: Knothe--Rosenblatt}\label{subsec:6.3}

\emph{Optimal-transport context, briefly.} Between two probability
measures on \(\mathbb{R}^d\), the \textbf{Knothe--Rosenblatt
rearrangement} (Rosenblatt 1952) is the unique transport map of
triangular form: each coordinate transformation conditions on
previously-transformed coordinates. It differs from the \textbf{Brenier
map} (Brenier 1991), which is the unique transport map that is the
gradient of a convex function (the \(L^2\)-optimal transport plan). KR
depends on coordinate ordering; Brenier does not. Both are bijections
between the two measures; KR is the natural fit for autoregressive
architectures because its triangular structure aligns with the network's
autoregressive masking.

\textbf{Theorem A-d (Multivariate uniqueness within the autoregressive
triangular class).} \emph{Let
\(X \mid \theta \sim p(\cdot \mid \theta)\) be a continuous model with
the property that, for each \(k\), the conditional CDF
\(F_k^{(\theta)}(\cdot \mid X_{<k}) := P(X_k \leq \cdot \mid \theta, X_{<k})\)
is \(C^1\) and strictly increasing on the support of
\(X_k \mid \theta, X_{<k}\). Within the class of autoregressive
triangular pivots satisfying (R1\(^\mathrm{auto}\)) and
(R2\(^\mathrm{auto}\)) with each \(r_k \in C^1\) in its
\(X_k\)-argument, the calibration manifold \(\mathcal{M}\) consists of a
unique function --- the \textbf{Knothe--Rosenblatt rearrangement}:} \[
r_k^\mathrm{KR}(\theta, X) = \Phi^{-1}\bigl(1 - F_k^{(\theta)}(X_k \mid X_{<k})\bigr).
\]

\emph{Proof.} Induct on \(k\).

\emph{Base case (\(k = 1\)).} Under (R2\(^\mathrm{auto}\)),
\(r_1: \Theta \times \mathcal{X}_1 \to \mathbb{R}\) depends on \(X\)
only through \(X_1\). It may depend on \(\theta\) generally, but by
(R1\(^\mathrm{auto}\)) only through \(\theta_1\). The full calibration
constraint requires
\(r_1(\theta, X_1) \mid \theta \sim \mathcal{N}(0,1)\). Apply Theorem A
in \(X_1\)-coordinates at each fixed \(\theta_1\) (with other
\(\theta\)-components irrelevant by (R1\(^\mathrm{auto}\))): the unique
monotone-in-\(X_1\) pushforward of the marginal \(X_1 \mid \theta_1\) to
\(\mathcal{N}(0,1)\) is \(\Phi^{-1}(1 - F_1^{(\theta_1)}(X_1))\).
(R1\(^\mathrm{auto}\)) fixes the sign.

\emph{Inductive step (\(k > 1\)).} Assume \(r_1, \ldots, r_{k-1}\) are
uniquely identified. Consider \(r_k(\theta, X_{\le k})\) at fixed
\((\theta, X_{<k})\). The calibration constraint
\(r(\theta; X) \mid \theta \sim \mathcal{N}(0, I_d)\) implies that the
components are jointly standard normal; since \(r_1, \ldots, r_{k-1}\)
are deterministic functions of \((\theta, X_{<k})\), conditioning on
them (given \(\theta\)) is equivalent to conditioning on \(X_{<k}\).
Hence
\(r_k(\theta, X_{\le k}) \mid \theta, X_{<k} \sim \mathcal{N}(0,1)\). By
(R2\(^\mathrm{auto}\)), as a function of \(X_k\) this map is monotone;
by Theorem A applied in \(X_k\)-coordinates with conditional law
\(F_k^{(\theta)}(\cdot \mid X_{<k})\) playing the role of the Gaussian
source, it is uniquely
\(\Phi^{-1}(1 - F_k^{(\theta)}(X_k \mid X_{<k}))\).
(R1\(^\mathrm{auto}\)) fixes the sign of the dependence on \(\theta_k\).
\(\blacksquare\)

\textbf{Corollary (multivariate location-Gaussian).} \emph{For
\(X \mid \theta \sim \mathcal{N}(\theta, \Sigma)\) with \(\Sigma\)
positive definite, let \(LL^T = \Sigma\) be the \textbf{Cholesky
factorization} (\(L\) lower-triangular with positive diagonal). Then the
KR pivot is} \[
r^\mathrm{KR}(\theta, X) = L^{-1}(\theta - X).
\]

\emph{Proof.} The conditional \(X_k \mid \theta, X_{<k}\) is Gaussian
with mean
\(\mu_k = \theta_k + \sum_{j < k} \alpha_{kj}(X_j - \theta_j)\) and
variance \(\sigma_k^2 = L_{kk}^2\), where the
\(\alpha_{kj}, \sigma_k^2\) are the Cholesky regression coefficients.
Direct computation:
\(F_k^{(\theta)}(X_k \mid X_{<k}) = \Phi((X_k - \mu_k)/\sigma_k)\), so
\[
r_k^\mathrm{KR} = \Phi^{-1}(1 - F_k) = \frac{\mu_k - X_k}{\sigma_k} = \frac{\theta_k - X_k + \sum_{j < k} \alpha_{kj}(X_j - \theta_j)}{L_{kk}} = [L^{-1}(\theta - X)]_k.
\] The verification on the 2D correlated experiment
(\(\Sigma = \bigl(\!\begin{smallmatrix} 1 & 0.5 \\ 0.5 & 1\end{smallmatrix}\!\bigr)\),
giving
\(L^{-1} = \bigl(\!\begin{smallmatrix} 1 & 0 \\ -0.577 & 1.155\end{smallmatrix}\!\bigr)\),
hence
\(r_2^\mathrm{KR} = 1.155(\theta_2 - X_2) - 0.577(\theta_1 - X_1)\)) is
reported in \S\ref{subsec:8.3}. \(\blacksquare\)

\textbf{Remarks.}

\emph{1. KR depends on coordinate ordering.} Different orderings of the
coordinates of \(X\) (and correspondingly of \(\theta\) at each
autoregressive level) yield different KR pivots, all valid CD pivots,
all in their respective calibration manifolds. The architectural choice
of autoregressive order in the network implicitly selects one. For
models with a natural ordering (time series; spatial data with a fixed
sweep), use the data-defined order; for unordered data, any ordering is
correct, with different choices giving different but inferentially
equivalent pivots.

\emph{2. ``Calibration manifold is a singleton'' vs.~empirical pointwise
residuals.} The theorem says the function \(r\) is uniquely determined
within the class. The architectural decomposition \(r_k = a_k - b_k\)
(additive form) is \emph{not} uniquely determined --- one can shift mass
between \(a_k\) and \(b_k\) via any context-dependent function
\(h(\theta_{<k}, X_{<k})\), since \(a_k\) and \(b_k\) share that
context. Only the difference \(r_k = a_k - b_k\) is identified. The
pointwise residuals seen in the 2D correlated experiment (\S\ref{subsec:8.3}) reflect
finite-network approximation error around a unique target
\(r^\mathrm{KR}\), not non-uniqueness of the population-level solution.

\emph{3. The role of (R2\(^\mathrm{auto}\)) in multivariate.} Without
monotonicity in \(X_k\) at fixed \((\theta, X_{<k})\), the 1D inductive
step fails: \(X_k \mapsto r_k\) could be non-bijective, and the
conditional pushforward density would no longer be a monotone
rearrangement. This is the multivariate analogue of the
autograd-Jacobian failure mode described in \S\ref{subsec:8.4} --- the same
architectural prescription that fixes the 1D non-additive case is what
makes the inductive proof go through in \(d > 1\).

\subsection{Inference: direct pivot inversion}\label{subsec:6.4}

At an observed \(X_\mathrm{obs}\), the level-\(\alpha\) confidence set
is \[
\mathcal{C}_\alpha(X_\mathrm{obs}) := \bigl\{\theta : \|r(\theta; X_\mathrm{obs})\|^2 \leq \chi^2_{d, \alpha}\bigr\}.
\] Under (C1), \(\|r(\theta_0; X)\|^2 \mid \theta_0 \sim \chi^2_d\), so
\(\mathcal{C}_\alpha\) has exact frequentist coverage: \[
\mathbb{P}_{X \sim P_{\theta_0}}\bigl(\theta_0 \in \mathcal{C}_\alpha(X)\bigr) = \alpha.
\] The set \(\mathcal{C}_\alpha\) is found by numerical root-finding in
\(\theta\): since \(r\) is monotone in \(\theta\) (autoregressively),
each ``slice'' \(\theta_k \mid \theta_{<k}\) is a 1D root-finding
problem.

\textbf{Worked example (1D location-normal).} With
\(r^*(\theta; X) = \theta - X\),
\(\|r^*(\theta; X_\mathrm{obs})\|^2 = (\theta - X_\mathrm{obs})^2 \le \chi^2_{1, \alpha} = z_\alpha^2\)
gives
\(\mathcal{C}_\alpha = [X_\mathrm{obs} - z_\alpha, X_\mathrm{obs} + z_\alpha]\)
--- the standard \(z\)-interval, as expected.

\medskip\hrule\medskip

\section*{Part V — Empirical validation}
\addcontentsline{toc}{section}{Part V — Empirical validation}

\section{Architecture and training}\label{sec:7}

\subsection{UMNN building blocks}\label{subsec:7.1}

\textbf{Definition (UMNN; Wehenkel \& Louppe 2019).} A
\emph{Unconstrained Monotonic Neural Network} parameterizes a scalar
function \(g(z; c)\) that is strictly monotone-increasing in \(z\) at
fixed context \(c\) by writing \[
g(z; c) = \mathrm{bias}(c) + \int_{z_0}^z \mathrm{softplus}\bigl(\mathrm{MLP}(t, c)\bigr)\, dt,
\] where \(\mathrm{MLP}(\cdot, c)\) is an unconstrained neural network.
The integrand is positive everywhere (softplus is positive), so \(g\) is
strictly monotone in \(z\) by construction. The derivative
\(\partial_z g(z; c) = \mathrm{softplus}(\mathrm{MLP}(z, c))\) is
closed-form --- the Jacobian factor needed for NF-MLE is available
without recursive autograd.

We use Gauss--Legendre quadrature (\(n = 12\) nodes by default) to
evaluate the integral.

Hidden width 32 throughout. Final-layer weights initialized to zero so
the bias dominates at init and the model starts near a near-identity
transformation. We multiply the integrated function by a learnable
positive scalar \(\alpha\) (initialized at \(0.5 / \ln 2 \approx 0.72\),
giving init effective slope \(0.5\), deliberately wrong) so optimization
has to find the right scale.

\subsection{Training schedule}\label{subsec:7.2}

\begin{itemize}
\tightlist
\item
  Batch size 512, Adam at learning rate \(3 \times 10^{-3}\), gradient
  clipping at norm 5.
\item
  3,000--4,000 steps total. On CPU with 4 threads, this runs in 2--3
  minutes per experiment.
\item
  Single-stage training. No annealing schedule (NF-MLE doesn't need
  one).
\end{itemize}

\subsection{Diagnostics}\label{subsec:7.3}

For each trained model:

\begin{enumerate}
\def\labelenumi{\arabic{enumi}.}
\tightlist
\item
  \textbf{Pivot RMSE vs \(r^*\)} on a dense grid, where \(r^*\) is
  closed-form. Sanity check that the model recovered the analytical
  answer.
\item
  \textbf{Marginal PIT.} \(\Phi(r_k(\theta_i, X_i))\) on
  training-distribution samples should be \(U(0,1)\).
  Kolmogorov--Smirnov (KS) test.
\item
  \textbf{Conditional PIT.} At fixed \(\theta_0\),
  \(r_k(\theta_0, X) \mid \theta_0\) should be \(\mathcal{N}(0,1)\). KS
  test at \(\geq 5\) \(\theta_0\) values.
\item
  \textbf{Joint Mahalanobis (\(d > 1\)).}
  \(\|r(\theta_0, X)\|^2 \mid \theta_0\) should be \(\chi^2_d\). KS
  test.
\item
  \textbf{Coverage.}
  \(\mathbb{P}_{X \sim P_{\theta_0}}(\theta_0 \in \mathcal{C}_\alpha(X))\)
  across \(\theta_0\) and confidence levels
  \(\alpha \in \{0.5, 0.68, 0.9, 0.95\}\). This is the inferentially
  relevant test.
\end{enumerate}

\textbf{Statistical noise floor for KS.} Under the null,
\(\sqrt{N} \cdot \mathrm{KS}\) has the Kolmogorov limiting distribution
with \(\Pr(\sqrt{N}\cdot\mathrm{KS} > 1.628) = 0.01\). At \(N = 5000\),
this gives a noise floor of
\(\mathrm{KS} \approx 1.628/\sqrt{5000} \approx 0.023\) for the
0.01-level threshold; at \(N = 6000\), the floor is \(\approx 0.021\).
We refer to this floor when interpreting ``borderline'' results in \S\ref{sec:8}.

Related diagnostic frameworks include simulation-based calibration (SBC;
Talts et al.~2018), sampling-based accuracy testing (TARP; Lemos et
al.~2023), and the LF2I conditional-coverage diagnostic (Dalmasso et
al.~2024). Our coverage diagnostic (item 5) is closest in spirit to the
LF2I local-coverage check.

\medskip\hrule\medskip

\section{Four experiments}\label{sec:8}

Each experiment is supporting evidence for a specific theoretical claim
from Parts III--IV. They are not validation in the sense of ``checking
that the method works'' --- the theory already establishes correctness
in the population limit. They validate that (i) finite-sample
optimization actually reaches the population optimum at practical
training budgets, and (ii) the architectural prescriptions matter as the
theory says they should.

\begin{longtable}[]{@{}
  >{\raggedright\arraybackslash}p{(\columnwidth - 2\tabcolsep) * \real{0.5000}}
  >{\raggedright\arraybackslash}p{(\columnwidth - 2\tabcolsep) * \real{0.5000}}@{}}
\toprule\noalign{}
\begin{minipage}[b]{\linewidth}\raggedright
Experiment
\end{minipage} & \begin{minipage}[b]{\linewidth}\raggedright
Supports
\end{minipage} \\
\midrule\noalign{}
\endhead
\bottomrule\noalign{}
\endlastfoot
\S\ref{subsec:8.1} 1D location-normal & Theorem A (1D uniqueness in
\(\mathcal{F}_{C^1}\)) \\
\S\ref{subsec:8.2} 2D location, \(\Sigma = I\) & Triangular flow construction
(\S\ref{subsec:6.1}--6.2) on a clean diagonal case \\
\S\ref{subsec:8.3} 2D location, \(\Sigma\) correlated & Theorem A-d / KR uniqueness
(\S\ref{subsec:6.3}), via recovered Jacobian \(= L^{-1}\) \\
\S\ref{subsec:8.4} Exponential rate & Architectural prescription (R2): the
doubly-monotone construction succeeds, the autograd-Jacobian ablation
fails \\
\end{longtable}

Each experiment trains in 2--4 minutes on CPU.

\subsection{1D location-normal}\label{subsec:8.1}

\textbf{Setup.} \(X \sim \mathcal{N}(\theta, 1)\),
\(\theta \sim U[-7, 7]\). Truth \(r^*(\theta, X) = \theta - X\).

\textbf{Architecture.} Additive flow
\(r(\theta, X) = \alpha_a \cdot a(\theta) - \alpha_b \cdot b(X)\) with
\(a, b\) scalar UMNNs.

\textbf{Results.}

\begin{longtable}[]{@{}
  >{\raggedright\arraybackslash}p{(\columnwidth - 4\tabcolsep) * \real{0.3333}}
  >{\raggedright\arraybackslash}p{(\columnwidth - 4\tabcolsep) * \real{0.3333}}
  >{\raggedright\arraybackslash}p{(\columnwidth - 4\tabcolsep) * \real{0.3333}}@{}}
\toprule\noalign{}
\begin{minipage}[b]{\linewidth}\raggedright
Diagnostic
\end{minipage} & \begin{minipage}[b]{\linewidth}\raggedright
Value
\end{minipage} & \begin{minipage}[b]{\linewidth}\raggedright
Status
\end{minipage} \\
\midrule\noalign{}
\endhead
\bottomrule\noalign{}
\endlastfoot
Pivot RMSE & 0.030 & clean (\textasciitilde0.3\% of dynamic range) \\
Slope at \(\theta = 0, X = 0\) & \(\partial_\theta r = 0.997\),
\(\partial_X r = -0.983\) & within 1--2\% of truth \\
Marginal PIT KS & 0.008 & PASS (well below noise floor) \\
Conditional PIT KS at 7 \(\theta_0\) values & bulk KS \(\approx 0.01\),
edge KS up to 0.05 & bulk PASS; edge effects at \(\rho\) boundary \\
Coverage error at all levels in interior & \(< 0.01\) & excellent \\
\end{longtable}

The 1D case is a sanity check. Both the additive architecture and a
(more flexible) joint UMNN architecture converge under NF-MLE, with the
additive form converging 20× faster and to lower RMSE because it exactly
contains \(r^*\). The conditional PIT shows degradation near the support
boundary of \(\rho\) --- a known finite-sample effect from undersampling
near the support edges, not a structural issue.

\subsection{2D location-Gaussian, \(\Sigma = I\)}\label{subsec:8.2}

\textbf{Setup.} \(X \mid \theta \sim \mathcal{N}(\theta, I_2)\),
\(\theta \sim U[-7, 7]^2\). Truth \(r^*(\theta, X) = \theta - X\)
component-wise.

\textbf{Architecture.} Triangular flow as in \S\ref{subsec:6.1} (additive form):
\[r_1(\theta, X) = a_1(\theta_1) - b_1(X_1), \qquad r_2(\theta, X) = a_2(\theta_2; \theta_1, X_1) - b_2(X_2; \theta_1, X_1).\]
\(a_1, b_1\) scalar UMNNs; \(a_2, b_2\) contextual UMNNs taking
\((\theta_1, X_1)\) as context. 5,004 parameters total.

\textbf{Results.}

\begin{longtable}[]{@{}
  >{\raggedright\arraybackslash}p{(\columnwidth - 4\tabcolsep) * \real{0.3333}}
  >{\raggedright\arraybackslash}p{(\columnwidth - 4\tabcolsep) * \real{0.3333}}
  >{\raggedright\arraybackslash}p{(\columnwidth - 4\tabcolsep) * \real{0.3333}}@{}}
\toprule\noalign{}
\begin{minipage}[b]{\linewidth}\raggedright
Diagnostic
\end{minipage} & \begin{minipage}[b]{\linewidth}\raggedright
Value
\end{minipage} & \begin{minipage}[b]{\linewidth}\raggedright
Status
\end{minipage} \\
\midrule\noalign{}
\endhead
\bottomrule\noalign{}
\endlastfoot
Pivot RMSE (total) & 0.044 & clean \\
Per-component: \(r_1\) & 0.014 & excellent \\
Per-component: \(r_2\) & 0.061 & clean (conditioning network learning
slower) \\
Marginal PIT KS: \(r_1, r_2\) & 0.007, 0.006 & both PASS
(\(p > 0.1\)) \\
Joint Mahalanobis vs \(\chi^2_2\) KS & 0.011 & clean \\
Conditional PIT at 5 \(\theta_0\): \(r_1\) & KS 0.011--0.018, all PASS &
clean \\
Conditional PIT at 5 \(\theta_0\): \(r_2\) & KS 0.007--0.044, 3/5
borderline at \(\alpha=0.01\) & KS values near the \(N=5000\) noise
floor (0.023) \\
Conditional Mahalanobis at 5 \(\theta_0\) & KS 0.008--0.016, all PASS &
clean \\
Coverage error across \(\theta_0\), all \(\alpha\) & \(< 0.01\) &
excellent \\
\end{longtable}

The clean result. Triangular construction scales naturally, and the
joint Mahalanobis test (which \(r_1, r_2\) being individually calibrated
does \emph{not} automatically imply) passes. The ``borderline''
conditional PIT for \(r_2\) reflects KS values within \textasciitilde2×
the noise floor at \(N = 5000\); the more inferentially relevant joint
Mahalanobis and coverage tests are clean.

\subsection{2D location-Gaussian, \(\Sigma\) correlated}\label{subsec:8.3}

\textbf{Setup.}
\(\Sigma = \begin{pmatrix} 1 & 0.5 \\ 0.5 & 1 \end{pmatrix}\). Truth
\(r^*(\theta, X) = L^{-1}(\theta - X)\) where \(LL^T = \Sigma\).
Explicitly:
\[r_1^* = \theta_1 - X_1, \qquad r_2^* = -0.577(\theta_1 - X_1) + 1.155(\theta_2 - X_2).\]
The conditioning network for \(r_2\) now must learn nontrivial
cross-coupling on \((\theta_1, X_1)\), with slope \(-0.577\).

\textbf{Results.}

\begin{longtable}[]{@{}
  >{\raggedright\arraybackslash}p{(\columnwidth - 4\tabcolsep) * \real{0.3333}}
  >{\raggedright\arraybackslash}p{(\columnwidth - 4\tabcolsep) * \real{0.3333}}
  >{\raggedright\arraybackslash}p{(\columnwidth - 4\tabcolsep) * \real{0.3333}}@{}}
\toprule\noalign{}
\begin{minipage}[b]{\linewidth}\raggedright
Diagnostic
\end{minipage} & \begin{minipage}[b]{\linewidth}\raggedright
Value
\end{minipage} & \begin{minipage}[b]{\linewidth}\raggedright
Status
\end{minipage} \\
\midrule\noalign{}
\endhead
\bottomrule\noalign{}
\endlastfoot
Pivot RMSE (total) & 0.27 & ``large'' (but see commentary) \\
Marginal PIT KS: \(r_1, r_2\) & 0.010, 0.003 & both PASS \\
Joint Mahalanobis vs \(\chi^2_2\) KS & 0.012 & clean \\
Coverage error across all \(\theta_0\), \(\alpha\) & \(< 0.015\) &
excellent \\
Recovered \(\partial_\theta r\) at 200 random points &
\(\mathbb{E}[J_\theta] = \begin{pmatrix} 1.000 & 0 \\ -0.569 & 1.154\end{pmatrix} \pm \begin{pmatrix} 0.004 & 0 \\ 0.049 & 0.007\end{pmatrix}\)
& matches \(L^{-1}\) to 1--2\% \\
\end{longtable}

The ``high'' pivot RMSE of 0.27 is misleading: \(r_2\) has dynamic range
\(\sim 12\), so 0.27 is \(\sim 3\%\) relative. The recovered linear
Jacobian coefficients match \(L^{-1}\) to 1--2\% on average, with
residual variance \(\sim 0.05\) in the off-diagonal entries reflecting
finite-network approximation error around the \emph{unique} KR target
(Theorem A-d, \S\ref{subsec:6.3}). Coverage is exact to within 1.5\%.

\textbf{What this validates.} Theorem A-d says the calibration manifold
within the autoregressive triangular class is the singleton
Knothe--Rosenblatt rearrangement. For Gaussian models, KR equals
\(L^{-1}(\theta - X)\) (corollary in \S\ref{subsec:6.3}). The recovered average
Jacobian matches this exactly; the pointwise residuals are
finite-network noise around the population-level singleton, not evidence
of non-uniqueness. The proper diagnostic in multivariate is the
recovered Jacobian (matches truth to 1--2\%) and coverage (within
1.5\%), not the pointwise pivot RMSE (which conflates approximation
error with non-canonical alignment).

\subsection{Exponential rate: validating the non-additive case}\label{subsec:8.4}

\textbf{Setup.} \(X_i \sim \mathrm{Exp}(\theta)\) iid, \(n = 5\) samples
summed into \(T = \sum_i X_i\). \(\theta \sim U[0.3, 3.0]\). Truth
\(r^*(\theta, T) = \Phi^{-1}(F_{\chi^2_{10}}(2\theta T))\) ---
non-additive, with multiplicative \(\theta T\) interaction. This tests
whether the framework handles models outside the additive class of
\S\ref{subsec:8.1}--8.3.

\textbf{Doubly-monotone architecture (\S\ref{subsec:6.1}, form 2).} Both
monotonicities enforced architecturally, with \(X\)-dependence through
the sufficient statistic \(T\): \[
r(\theta, T) = b_\mathrm{umnn}(T) + \int_{\theta_\mathrm{ref}}^{\theta} \mathrm{softplus}\bigl(\alpha(t) + \beta_\mathrm{umnn}(T)\bigr)\, dt,
\] where \(b_\mathrm{umnn}, \beta_\mathrm{umnn}\) are
monotone-increasing UMNNs of \(T\). Then
\(\partial_\theta r = \mathrm{softplus}(\cdot) > 0\) (R1) and
\(\partial_T r = b'_\mathrm{umnn}(T) + \beta'_\mathrm{umnn}(T) \int_{\theta_\mathrm{ref}}^\theta \sigma(\cdot)\, dt > 0\)
(R2), both in closed form.

\begin{longtable}[]{@{}
  >{\raggedright\arraybackslash}p{(\columnwidth - 4\tabcolsep) * \real{0.3333}}
  >{\raggedright\arraybackslash}p{(\columnwidth - 4\tabcolsep) * \real{0.3333}}
  >{\raggedright\arraybackslash}p{(\columnwidth - 4\tabcolsep) * \real{0.3333}}@{}}
\toprule\noalign{}
\begin{minipage}[b]{\linewidth}\raggedright
Diagnostic
\end{minipage} & \begin{minipage}[b]{\linewidth}\raggedright
Value
\end{minipage} & \begin{minipage}[b]{\linewidth}\raggedright
Status
\end{minipage} \\
\midrule\noalign{}
\endhead
\bottomrule\noalign{}
\endlastfoot
Pivot RMSE on natural support & 0.045 & clean \\
Final loss & 0.87 & matches truth's 0.88 to batch noise \\
Marginal PIT KS & 0.006 (\(p = 0.23\)) & PASS \\
Conditional PIT KS at 5 \(\theta_0\) & 0.016--0.028 & borderline
(\(N=6000\) noise floor \(\approx 0.021\)) \\
Coverage error across \(\theta_0, \alpha\) & \(< 0.015\) & excellent \\
\end{longtable}

The doubly-monotone construction recovers \(r^*\) to within batch noise,
demonstrating that the framework extends beyond the additive class.

\textbf{Architectural ablation: why (R2) must be enforced
architecturally.} Replacing the doubly-monotone construction with a
joint UMNN in \(\theta\) that uses autograd to compute \(\partial_T r\)
--- i.e., no architectural enforcement of monotone-in-\(T\) --- yields
RMSE 1.56, conditional PIT KS 0.045--0.080 (all FAIL relative to the
noise floor), and final loss 0.56. The loss value 0.56 is \emph{below}
the analytical truth's loss of 0.88, which would be impossible if the
ablation model lived in a valid (bijective-in-\(X\)) class. The
mechanism: when \(T \mapsto r\) is non-bijective, \(\hat p_r\) is not
normalized in \(T\) at fixed \(\theta\), so \(-\log \hat p_r\) has no
lower bound at the conditional entropy. Concretely, in regions where the
autograd \(|\partial r / \partial T|\) is locally inflated by a fold,
\(\log|\partial r / \partial T|\) contributes a large positive value to
the surrogate log-density without \(\hat p_r\) being a valid density.
This is direct empirical evidence for (R2): architectural enforcement of
monotone-in-\(X\) is not a stylistic preference but a correctness
requirement for the NF-MLE objective.

\subsection{Synthesis}\label{subsec:8.5}

The four experiments collectively establish that the population-level
theory of Parts III--IV is reached at practical finite-sample budgets
(\(N \in [5000, 10000]\), training time 2--4 min on CPU). Mapping
experiments back to claims:

\begin{longtable}[]{@{}
  >{\raggedright\arraybackslash}p{(\columnwidth - 2\tabcolsep) * \real{0.5000}}
  >{\raggedright\arraybackslash}p{(\columnwidth - 2\tabcolsep) * \real{0.5000}}@{}}
\toprule\noalign{}
\begin{minipage}[b]{\linewidth}\raggedright
Theoretical claim
\end{minipage} & \begin{minipage}[b]{\linewidth}\raggedright
Empirical evidence
\end{minipage} \\
\midrule\noalign{}
\endhead
\bottomrule\noalign{}
\endlastfoot
Theorem A: \(r^* = \theta - X\) uniquely calibrates in 1D & \S\ref{subsec:8.1}: RMSE
0.030; slopes recovered to 1--2\%; coverage error \(< 0.01\) \\
Triangular flow construction yields a valid CD pivot in \(d > 1\)
(\S\ref{subsec:6.1}--6.2) & \S\ref{subsec:8.2}: joint Mahalanobis KS 0.011, coverage error
\(< 0.01\) \\
Theorem A-d: KR \(= L^{-1}(\theta - X)\) is the unique pivot in
triangular class for Gaussian (\S\ref{subsec:6.3} corollary) & \S\ref{subsec:8.3}: recovered
\(\mathbb{E}[J_\theta]\) matches \(L^{-1}\) to 1--2\%; coverage exact to
within 1.5\% \\
(R2) is a correctness requirement, not a stylistic preference & \S\ref{subsec:8.4}:
doubly-monotone construction recovers \(r^*\) (RMSE 0.045);
autograd-Jacobian ablation produces a loss \emph{below} truth and fails
calibration \\
\end{longtable}

Summary table:

\begin{longtable}[]{@{}
  >{\raggedright\arraybackslash}p{(\columnwidth - 8\tabcolsep) * \real{0.2000}}
  >{\raggedright\arraybackslash}p{(\columnwidth - 8\tabcolsep) * \real{0.2000}}
  >{\raggedright\arraybackslash}p{(\columnwidth - 8\tabcolsep) * \real{0.2000}}
  >{\raggedright\arraybackslash}p{(\columnwidth - 8\tabcolsep) * \real{0.2000}}
  >{\raggedright\arraybackslash}p{(\columnwidth - 8\tabcolsep) * \real{0.2000}}@{}}
\toprule\noalign{}
\begin{minipage}[b]{\linewidth}\raggedright
Experiment
\end{minipage} & \begin{minipage}[b]{\linewidth}\raggedright
Pivot RMSE
\end{minipage} & \begin{minipage}[b]{\linewidth}\raggedright
Marg KS
\end{minipage} & \begin{minipage}[b]{\linewidth}\raggedright
Coverage error
\end{minipage} & \begin{minipage}[b]{\linewidth}\raggedright
Theoretical claim supported
\end{minipage} \\
\midrule\noalign{}
\endhead
\bottomrule\noalign{}
\endlastfoot
1D location-normal & 0.030 & 0.008 & \(< 0.01\) & Theorem A \\
2D loc, \(\Sigma = I\) & 0.044 & 0.007, 0.006 & \(< 0.01\) & \S\ref{subsec:6.1}
construction \\
2D loc, \(\Sigma\) corr. & 0.27 (3\% rel.) & 0.010, 0.003 & \(< 0.015\)
& Theorem A-d (Jacobian \(= L^{-1}\)) \\
Exp-rate (doubly-mono) & 0.045 & 0.006 & \(< 0.015\) & (R2)
sufficiency \\
Exp-rate (autograd Jac, ablation) & 1.56 & 0.024 & up to 3.4\% & (R2)
necessity \\
\end{longtable}

The architectural prescription distilled across all four cases:

\begin{enumerate}
\def\labelenumi{\arabic{enumi}.}
\tightlist
\item
  \textbf{NF-MLE loss} (\S\ref{sec:3}, Lemma 4.3).
\item
  \textbf{Monotone in \(\theta\), architecturally} (R1) --- UMNN
  integration in \(\theta\).
\item
  \textbf{Monotone in \(X\), architecturally} (R2) --- additive
  structure with monotone-in-\(X\) heads, triangular flow in \(d > 1\),
  or doubly-monotone UMNN for non-additive \(X\)-dependence.
\end{enumerate}

\medskip\hrule\medskip

\section*{Part VI — Practice and position}
\addcontentsline{toc}{section}{Part VI — Practice and position}

\section{Implementation recipe}\label{sec:9}

A practitioner's checklist for applying CD-SBI to a new model:

\textbf{1. Identify the structure of the truth \(r^*\).} Is it additive
in \((\theta, X)\)? Multiplicative? Does the sufficient statistic factor
through? This tells you which architecture to use:

\begin{longtable}[]{@{}
  >{\raggedright\arraybackslash}p{(\columnwidth - 2\tabcolsep) * \real{0.5000}}
  >{\raggedright\arraybackslash}p{(\columnwidth - 2\tabcolsep) * \real{0.5000}}@{}}
\toprule\noalign{}
\begin{minipage}[b]{\linewidth}\raggedright
If truth is approximately
\end{minipage} & \begin{minipage}[b]{\linewidth}\raggedright
Use architecture
\end{minipage} \\
\midrule\noalign{}
\endhead
\bottomrule\noalign{}
\endlastfoot
Additive: \(r^* \approx a(\theta) - b(X)\) & Additive flow (\S\ref{subsec:6.1} form 1;
\S\ref{subsec:8.1}--8.2) \\
Multivariate with cross-coupling & Triangular flow, additive form (\S\ref{subsec:6.1}
form 1; \S\ref{subsec:8.3}) \\
Has multiplicative \(\theta\)--\(X\) interaction & Doubly-monotone
construction (\S\ref{subsec:6.1} form 2; \S\ref{subsec:8.4}) \\
No prior knowledge & Doubly-monotone construction with autoregressive
conditioning \\
\end{longtable}

\textbf{2. Enforce both monotonicities architecturally.} Never rely on
autograd to give you a valid Jacobian without architectural support. The
\S\ref{subsec:8.4} ablation is documentation of why.

\textbf{3. Use a proposal \(\rho\) with full support on the region of
interest.} The framework is robust to the choice of \(\rho\) in the
population limit, but finite-sample performance benefits from \(\rho\)
covering the region where inference will be performed. Edge effects
(\S\ref{subsec:8.1}) appear near the support boundary.

\textbf{4. Train with NF-MLE.} No annealing schedule needed. Batch size
256--1024; Adam at \(\sim 3 \times 10^{-3}\); 3,000--10,000 steps
depending on problem complexity.

\textbf{5. Validate.} Run all five diagnostics in \S\ref{subsec:7.3}. The critical one
is coverage: if coverage matches nominal at multiple \(\theta_0\) and
multiple \(\alpha\) levels, the pivot is valid. If only marginal PIT
passes but conditional fails, you have a Hermans-style problem and need
to revisit the architecture.

\textbf{6. At test time, invert.} Given observed \(X_\mathrm{obs}\),
compute
\(\mathcal{C}_\alpha = \{\theta : \|r(\theta, X_\mathrm{obs})\|^2 \le \chi^2_{d, \alpha}\}\)
by 1D root-finding (autoregressively, in \(d > 1\)).

\medskip\hrule\medskip

\section{Position in the SBI literature}\label{sec:10}

\begin{longtable}[]{@{}
  >{\raggedright\arraybackslash}p{(\columnwidth - 6\tabcolsep) * \real{0.2500}}
  >{\raggedright\arraybackslash}p{(\columnwidth - 6\tabcolsep) * \real{0.2500}}
  >{\raggedright\arraybackslash}p{(\columnwidth - 6\tabcolsep) * \real{0.2500}}
  >{\raggedright\arraybackslash}p{(\columnwidth - 6\tabcolsep) * \real{0.2500}}@{}}
\toprule\noalign{}
\begin{minipage}[b]{\linewidth}\raggedright
Approach
\end{minipage} & \begin{minipage}[b]{\linewidth}\raggedright
Target
\end{minipage} & \begin{minipage}[b]{\linewidth}\raggedright
Calibration
\end{minipage} & \begin{minipage}[b]{\linewidth}\raggedright
Reference
\end{minipage} \\
\midrule\noalign{}
\endhead
\bottomrule\noalign{}
\endlastfoot
NPE / SNPE & Bayesian posterior & None built-in & Papamakarios \& Murray
2016; Greenberg et al.~2019 \\
NLE / SNL & Likelihood & None built-in & Papamakarios et al.~2019 \\
NRE / SNRE & Likelihood ratio & None built-in & Hermans et al.~2020;
Miller et al.~2022 \\
Balanced NRE & Posterior & Conservative regularizer & Delaunoy et
al.~2022 \\
Calibrated NPE & Posterior & Differentiable coverage & Falkiewicz et
al.~2023 \\
LF2I (ACORE, BFF) & Confidence sets & Neyman construction on learned
test stat & Dalmasso et al.~2024 \\
WALDO & Confidence sets & Neyman inversion on Wald-style stat &
Masserano et al.~2023 \\
Box CD & Confidence sets & Depth-based, hyper-rectangles & Bortolato \&
Ventura 2025 \\
Variational SBI w/ coverage & Posterior & Conformal post-hoc & Patel et
al.~2023 \\
\textbf{CD-SBI (this work)} & \textbf{Confidence distribution} &
\textbf{Pointwise, by construction} & --- \\
\end{longtable}

The closest cousins are LF2I and WALDO (which target test-statistic CDFs
to construct confidence sets via Neyman inversion) and SNL (which uses
the same NF-MLE loss but without the monotone-in-\(\theta\)
architectural constraint, and without the frequentist inference
machinery).

CD-SBI's specific contribution is the architectural recipe (monotone in
\(\theta\) + monotone in \(X\) + triangular structure for multivariate)
that converts the standard SNL training objective into a procedure that
produces a calibrated confidence distribution. The framework inherits
all of SNL's tooling and scalability while providing pointwise
frequentist coverage by construction. Unlike LF2I/WALDO, which require a
separately-trained calibration step (the ``critical-values branch''),
CD-SBI is single-stage: the same network output gives both the CD and
the confidence set, with the calibration enforced in the population
limit by the architecture + loss combination rather than by a hold-out
calibration sample.

The framework also relates to the implicit-likelihood-ratio tradition
(Cranmer et al.~2015): the optimal pivot in the location-normal model is
the standardized log-likelihood gradient, and in regular exponential
families more generally relates to the signed-root log-likelihood ratio
(Schweder \& Hjort 2016, ch.~5).

\medskip\hrule\medskip

\section*{Part VII — Open questions}
\addcontentsline{toc}{section}{Part VII — Open questions}

\section{Open problems}\label{sec:11}

\subsection{Uniqueness beyond the autoregressive triangular class}\label{subsec:11.1}
Theorem A-d (\S\ref{subsec:6.3}) establishes uniqueness within autoregressive
triangular flows: the calibration manifold is the singleton
Knothe--Rosenblatt rearrangement. But the \emph{full} calibration
manifold over more flexible architectures (e.g., dense flows,
non-autoregressive coupling layers) is generically not a singleton ---
orthogonal rotations of \(r\) preserve calibration since
\(\mathcal{N}(0, I_d)\) is rotation-invariant. Whether NF-MLE with
looser architectural constraints still converges to a canonical pivot,
or whether the practitioner must commit to a specific autoregressive
structure to guarantee uniqueness, is open. The practical implication,
supported by the 2D correlated experiment, is that \emph{some}
triangular ordering must be chosen architecturally.

\subsection{Ordering selection in KR}\label{subsec:11.2}
When there is no natural
ordering of the coordinates of \(X\), different choices of
autoregressive ordering give different valid pivots, all in
\(\mathcal{M}\), all yielding exact coverage. Whether some orderings
have better finite-sample optimization behavior --- easier-to-learn
conditionals, lower variance gradients --- is an empirical question
worth investigating. The optimal-transport literature
(Carlier--Galichon--Santambrogio 2010) connects KR maps with \(t \to 1\)
to Brenier maps, suggesting that a more isotropic transport target might
give better-behaved finite-sample training; this is unexplored in the
CD-SBI context.

\subsection{The \(C^0\) residual in 1D}\label{subsec:11.3}
In the \(C^1\) case, Lemma~4.2 derives strict monotonicity in \(X\) from
(R1) + calibration + regularity, so (R3) is automatic. Theorem~A*
requires (R3) as an extra assumption only in the \emph{Lipschitz} case,
where Lemma~4.2's blow-up argument fails (a Lipschitz V-shape gives only
a bounded jump in the pushforward density, not a divergence). Whether
(R3) follows from (R1) + Lipschitz + calibration alone --- without the
uniform lower-bound assumption \(|\partial_X r| \ge c > 0\) --- is open.
Lemma~4.2's generalization (\S\ref{subsec:4.2} remark) covers the
\(C^1\) case for arbitrary continuous positive source; in practice,
architectural enforcement of strict \(X\)-monotonicity (UMNN /
triangular flows) makes the question moot.

\subsection{Misspecification}\label{subsec:11.4}
When the simulator does not exactly
generate the data, the calibration target \(\mathcal{N}(0, I_d)\)
becomes wrong. Robust-SBI ideas (Schmon et al.~2020; Dellaporta et
al.~2022; Wehenkel et al.~2025) should integrate with CD-SBI but the
right framework is unclear. Specifically: under misspecification, the
population minimizer of NF-MLE is no longer the true CD pivot but rather
the moment-projection of the model density onto the architectural class.
What this object's coverage properties are --- and whether CD-SBI
inherits a degraded but still useful guarantee --- is an open question.

\subsection{Higher-dimensional scaling}\label{subsec:11.5}
The triangular flow scales
linearly in \(d\) for the Jacobian, but conditioning networks for
high-\(k\) components carry growing context dimension. For
\(d \gtrsim 10\), careful architecture design (sparse autoregressive
structure, attention-based conditioning) likely becomes important.
Empirical validation up to \(d = 10\) is a natural next milestone.

\subsection{Sequential / amortized variants}\label{subsec:11.6}
SBI has both amortized
(single network for all \(X_\mathrm{obs}\)) and sequential (refine
around a specific \(X_\mathrm{obs}\)) variants. CD-SBI as developed here
is amortized; a sequential variant would train \(\rho\) to concentrate
around \(\theta\) values consistent with \(X_\mathrm{obs}\). Whether the
calibration guarantee survives sequential refinement is non-trivial: the
implicit conditioning on the iterative procedure changes the
population-level distribution that the loss converges to.

\subsection{Discrete data with no scalar sufficient statistic}\label{subsec:11.7}
For
models where no scalar sufficient statistic exists (e.g., images,
sequences), the discrete extension of \S\ref{subsec:5.7} must be combined with feature
extraction. Whether the Lancaster randomization argument extends to
learned-feature settings is open.

\subsection{Efficient computation of \(\mathcal{C}_\alpha\) for
multivariate}\label{subsec:11.8}
Inversion via
\(\|r(\theta, X_\mathrm{obs})\|^2 \le \chi^2_{d, \alpha}\) requires
evaluating \(r\) on a grid or via root-finding. For visualization and
reporting, efficient confidence-set computation in \(d \geq 3\) ---
perhaps via direct sampling from the implicit set, or via a learned
set-boundary parameterization --- is worth developing.

\medskip\hrule\medskip

\section*{Status of claims}

\begin{longtable}[]{@{}
  >{\raggedright\arraybackslash}p{(\columnwidth - 4\tabcolsep) * \real{0.3333}}
  >{\raggedright\arraybackslash}p{(\columnwidth - 4\tabcolsep) * \real{0.3333}}
  >{\raggedright\arraybackslash}p{(\columnwidth - 4\tabcolsep) * \real{0.3333}}@{}}
\toprule\noalign{}
\begin{minipage}[b]{\linewidth}\raggedright
Claim
\end{minipage} & \begin{minipage}[b]{\linewidth}\raggedright
Status
\end{minipage} & \begin{minipage}[b]{\linewidth}\raggedright
Notes
\end{minipage} \\
\midrule\noalign{}
\endhead
\bottomrule\noalign{}
\endlastfoot
NF-MLE strictly proper for \(\mathcal{M}\) & Proven (\S\ref{subsec:3.2}; Lemma 4.3) &
KL non-negativity, conditional \\
Theorem A: uniqueness in 1D location-normal, \(C^1\) & Proven (\S\ref{subsec:4.3}) &
Three-lemma proof \\
Theorem A*: uniqueness in 1D, Lipschitz & Proven (\S\ref{subsec:4.4}) & Requires (R3)
as extra \\
Schweder--Hjort bridge & Proven (\S\ref{subsec:4.5}) &
\(r^* = \theta - X \Leftrightarrow C^*_\mathrm{SH}\) \\
Theorem C: lift to continuous one-parameter exponential families &
Proven (\S\ref{subsec:5.2}) & Same three-lemma structure in transformed coordinates \\
Theorem C*: discrete one-parameter, randomized & Proven (\S\ref{subsec:5.7}) &
Lex-order monotone rearrangement; requires (R4) \\
Theorem A-d: multivariate uniqueness within triangular class & Proven
(\S\ref{subsec:6.3}) & Calibration manifold \(=\) singleton KR; inductive reduction to
1D \\
Corollary: KR \(= L^{-1}(\theta - X)\) for multivariate
location-Gaussian & Proven (\S\ref{subsec:6.3}) & Verified empirically in \S\ref{subsec:8.3} \\
Empirical: 1D location-normal & Validated (\S\ref{subsec:8.1}) & RMSE 0.030, coverage
\(< 1\%\) error \\
Empirical: 2D location, \(\Sigma = I\) & Validated (\S\ref{subsec:8.2}) & RMSE 0.044,
joint Mahalanobis PASS \\
Empirical: 2D location, \(\Sigma\) correlated, recovers \(L^{-1}\) &
Validated (\S\ref{subsec:8.3}) & Jacobian within 1--2\% of KR; coverage exact \\
Empirical: exponential rate doubly-monotone & Validated (\S\ref{subsec:8.4}) & RMSE
0.045; autograd ablation produces direct evidence for (R2) \\
Uniqueness outside autoregressive triangular class & Open (\S\ref{subsec:11.1}) &
Rotational ambiguity in dense flows \\
KR ordering selection for unordered data & Open (\S\ref{subsec:11.2}) & Multiple valid
pivots; finite-sample comparison unclear \\
(R3) automatic for Lipschitz & Open (\S\ref{subsec:11.3}) & Architectural enforcement
sufficient in practice \\
Misspecification & Open (\S\ref{subsec:11.4}) & Not addressed \\
Higher-\(d\) scaling & Open (\S\ref{subsec:11.5}) & \(d \leq 2\) validated;
\(d \leq 10\) next milestone \\
Sequential refinement & Open (\S\ref{subsec:11.6}) & Theory unclear \\
\end{longtable}

\medskip\hrule\medskip

\section*{References}

\textbf{Confidence distributions and fiducial inference.}

\begin{itemize}
\tightlist
\item
  Cox, D. R. (1958). ``Some problems connected with statistical
  inference.'' \emph{Annals of Mathematical Statistics}, 29(2),
  357--372.
\item
  Efron, B. (1993). ``Bayes and likelihood calculations from confidence
  intervals.'' \emph{Biometrika}, 80(1), 3--26.
\item
  Efron, B. (1998). ``R. A. Fisher in the 21st century.''
  \emph{Statistical Science}, 13(2), 95--122.
\item
  Fisher, R. A. (1930). ``Inverse probability.'' \emph{Proceedings of
  the Cambridge Philosophical Society}, 26, 528--535.
\item
  Fraser, D. A. S. (2011). ``Is Bayes posterior just quick and dirty
  confidence?'' \emph{Statistical Science}, 26(3), 299--316.
\item
  Hwang, J. T. G., \& Yang, M.-C. (2001). ``An optimality theory for mid
  p-values in \(2 \times 2\) contingency tables.'' \emph{Statistica
  Sinica}, 11(3), 807--826.
\item
  Lancaster, H. O. (1961). ``Significance tests in discrete
  distributions.'' \emph{JASA}, 56(294), 223--234.
\item
  Schweder, T., \& Hjort, N. L. (2016). \emph{Confidence, Likelihood,
  Probability: Statistical Inference with Confidence Distributions}.
  Cambridge University Press.
\item
  Singh, K., Xie, M., \& Strawderman, W. E. (2007). ``Confidence
  distribution (CD) --- distribution estimator of a parameter.''
  \emph{IMS Lecture Notes --- Monograph Series}, 54, 132--150.
\item
  Xie, M., \& Singh, K. (2013). ``Confidence distribution, the
  frequentist distribution estimator of a parameter: A review.''
  \emph{International Statistical Review}, 81(1), 3--39.
\end{itemize}

\textbf{Simulation-based inference.}

\begin{itemize}
\tightlist
\item
  Bortolato, E., \& Ventura, L. (2025). ``Box Confidence Depth:
  Simulation-Based Inference with Hyper-Rectangles.'' arXiv:2502.11072.
\item
  Cranmer, K., Brehmer, J., \& Louppe, G. (2020). ``The frontier of
  simulation-based inference.'' \emph{PNAS}, 117(48), 30055--30062.
\item
  Cranmer, K., Pavez, J., \& Louppe, G. (2015). ``Approximating
  likelihood ratios with calibrated discriminative classifiers.''
  arXiv:1506.02169.
\item
  Dalmasso, N., Masserano, L., Zhao, D., Izbicki, R., \& Lee, A. B.
  (2024). ``Likelihood-free frequentist inference: bridging classical
  statistics and machine learning for reliable simulator-based
  inference.'' \emph{Electronic Journal of Statistics}, 18(2),
  5045--5090. arXiv:2107.03920.
\item
  Delaunoy, A., Hermans, J., Rozet, F., Wehenkel, A., \& Louppe, G.
  (2022). ``Towards reliable simulation-based inference with balanced
  neural ratio estimation.'' \emph{NeurIPS 35}.
\item
  Dellaporta, C., Knoblauch, J., Damoulas, T., \& Briol, F.-X. (2022).
  ``Robust Bayesian inference for simulator-based models via the MMD
  posterior bootstrap.'' \emph{AISTATS}.
\item
  Falkiewicz, M., Takeishi, N., Shekhzadeh, I., Wehenkel, A., Delaunoy,
  A., Louppe, G., \& Kalousis, A. (2023). ``Calibrating neural
  simulation-based inference with differentiable coverage probability.''
  \emph{NeurIPS 36}. arXiv:2310.13402.
\item
  Greenberg, D., Nonnenmacher, M., \& Macke, J. (2019). ``Automatic
  posterior transformation for likelihood-free inference.'' \emph{ICML}.
\item
  Hermans, J., Begy, V., \& Louppe, G. (2020). ``Likelihood-free MCMC
  with amortized approximate ratio estimators.'' \emph{ICML}.
\item
  Hermans, J., Delaunoy, A., Rozet, F., Wehenkel, A., Begy, V., \&
  Louppe, G. (2022). ``A trust crisis in simulation-based inference?
  Your posterior approximations can be unfaithful.'' \emph{Transactions
  on Machine Learning Research}. arXiv:2110.06581.
\item
  Lemos, P., Coogan, A., Hezaveh, Y., \& Perreault-Levasseur, L. (2023).
  ``Sampling-based accuracy testing of posterior estimators for general
  inference (TARP).'' \emph{ICML}.
\item
  Masserano, L., Dorigo, T., Izbicki, R., Kuusela, M., \& Lee, A. B.
  (2023). ``Simulator-based inference with WALDO: Confidence regions by
  leveraging prediction algorithms and posterior estimators for inverse
  problems.'' \emph{AISTATS}. arXiv:2205.15680.
\item
  Miller, B. K., Weniger, C., \& Forré, P. (2022). ``Contrastive neural
  ratio estimation.'' \emph{NeurIPS 35}.
\item
  Papamakarios, G., \& Murray, I. (2016). ``Fast \(\epsilon\)-free
  inference of simulation models with Bayesian conditional density
  estimation.'' \emph{NeurIPS}.
\item
  Papamakarios, G., Sterratt, D., \& Murray, I. (2019). ``Sequential
  neural likelihood.'' \emph{AISTATS}.
\item
  Patel, Y., McNamara, D., Loper, J., Regier, J., \& Tewari, A. (2023).
  ``Variational inference with coverage guarantees in simulation-based
  inference.'' arXiv:2305.14275.
\item
  Schmon, S. M., Cannon, P. W., \& Knoblauch, J. (2020). ``Generalized
  posteriors in approximate Bayesian computation.'' arXiv:2011.08644.
\item
  Talts, S., Betancourt, M., Simpson, D., Vehtari, A., \& Gelman, A.
  (2018). ``Validating Bayesian inference algorithms with
  simulation-based calibration.'' arXiv:1804.06788.
\item
  Wehenkel, A., Gamella, J. L., Sener, O., Behrmann, J., Sapiro, G.,
  Jacobsen, J.-H., \& Cuturi, M. (2025). ``Addressing misspecification
  in simulation-based inference through data-driven calibration.''
  \emph{ICML}. arXiv:2405.08719.
\end{itemize}

\textbf{Normalizing flows, monotone networks, optimal transport.}

\begin{itemize}
\tightlist
\item
  Brenier, Y. (1991). ``Polar factorization and monotone rearrangement
  of vector-valued functions.'' \emph{Comm. Pure Appl. Math.}, 44(4),
  375--417.
\item
  Carlier, G., Galichon, A., \& Santambrogio, F. (2010). ``From Knothe's
  transport to Brenier's map and a continuation method for optimal
  transport.'' \emph{SIAM J. Math. Anal.}, 41(6), 2554--2576.
\item
  Durkan, C., Bekasov, A., Murray, I., \& Papamakarios, G. (2019).
  ``Neural spline flows.'' \emph{NeurIPS 32}.
\item
  Papamakarios, G., Nalisnick, E., Rezende, D. J., Mohamed, S., \&
  Lakshminarayanan, B. (2021). ``Normalizing flows for probabilistic
  modeling and inference.'' \emph{JMLR}, 22, 1--64.
\item
  Rosenblatt, M. (1952). ``Remarks on a multivariate transformation.''
  \emph{Annals of Mathematical Statistics}, 23(3), 470--472.
\item
  Wehenkel, A., \& Louppe, G. (2019). ``Unconstrained monotonic neural
  networks.'' \emph{NeurIPS 32}. arXiv:1908.05164.
\end{itemize}

\textbf{Kernel methods and divergences.}

\begin{itemize}
\tightlist
\item
  Gneiting, T., \& Raftery, A. E. (2007). ``Strictly proper scoring
  rules, prediction, and estimation.'' \emph{JASA}, 102(477), 359--378.
\item
  Gretton, A., Bousquet, O., Smola, A., \& Schölkopf, B. (2005).
  ``Measuring statistical dependence with Hilbert--Schmidt norms.''
  \emph{Algorithmic Learning Theory}, 63--77.
\item
  Lyons, R. (2013). ``Distance covariance in metric spaces.''
  \emph{Annals of Probability}, 41(5), 3284--3305.
\item
  Sriperumbudur, B., Fukumizu, K., \& Lanckriet, G. (2011).
  ``Universality, characteristic kernels and RKHS embedding of
  measures.'' \emph{JMLR}, 12, 2389--2410.
\end{itemize}
\end{document}
